\pdfminorversion=7
\pdfobjcompresslevel=0
\documentclass[11pt, sans, a4paper]{moderncv}
\usepackage[utf8]{inputenc}
\usepackage[T1]{fontenc}
\usepackage{lmodern}
\usepackage{float}
\usepackage{ragged2e}
\usepackage{amsmath,amssymb}
\usepackage{booktabs}
\usepackage{array}
\usepackage{enumitem}
\usepackage{microtype}
\usepackage{adjustbox}
\usepackage{caption}
\usepackage[most]{tcolorbox}
\usepackage{hyperref}
\usepackage{catchfilebetweentags}
\usepackage{xr}
\externaldocument{manuscript_marked}

% Citation and reference bridge from manuscript_marked.aux
\makeatletter
\def\extractfirst#1#2#3#4{#1}
\def\bibcite#1#2{\@namedef{b@#1}{\extractfirst#2}}
\let\orig@newlabel\newlabel
\def\newlabel#1#2{}
\InputIfFileExists{manuscript_marked.aux}{}{}
\let\newlabel\orig@newlabel
% Suppress \label re-definitions from ExecuteMetaData-extracted equation blocks
% (labels are already loaded from manuscript_marked.aux; redefining them causes
% "multiply defined" warnings with no benefit since this letter has no \ref calls
% for equation numbers).
\AtBeginDocument{\let\label\@gobble}

% Fallback definitions for Overleaf (when manuscript_marked.aux is not generated in the same build)
\providecommand{\fallbackref}[2]{\@ifundefined{r@#1}{\@namedef{r@#1}{{#2}{}}}{}}
\providecommand{\fallbackcite}[2]{\@ifundefined{b@#1}{\@namedef{b@#1}{#2}}{}}

\fallbackref{eq:TM_binding}{1}
\fallbackref{eq:relative_stability}{2}
\fallbackref{fig:Fig1}{1}
\fallbackref{fig:Fig2}{2}
\fallbackref{fig:Fig3}{3}
\fallbackref{fig:Fig4}{4}
\fallbackref{fig:Fig5}{5}
\fallbackref{fig:Fig6}{6}
\fallbackref{fig:Fig7}{7}
\fallbackref{fig:Fig8}{8}

\fallbackcite{ASS2026_166886}{66}
\fallbackcite{Abdel-Hafiez2025}{24}
\fallbackcite{ApplSurfSci2025_161321}{59}
\fallbackcite{Gao2018}{36}
\fallbackcite{Ipaves2025}{7}
\fallbackcite{Liao2022}{6}
\fallbackcite{Luo2020}{35}
\fallbackcite{Mastrippolito2021}{25}
\fallbackcite{Moraes2026}{42}
\fallbackcite{NORSKOV2004}{43}
\fallbackcite{Norskov_2005}{41}
\fallbackcite{PCCP2018_15863}{60}
\fallbackcite{Vacuum2020_109438}{62}
\fallbackcite{Vacuum2020_109742}{61}
\fallbackcite{Webster2018}{21}
\fallbackcite{Zala2024}{40}
\makeatother

\moderncvstyle{classic}
\moderncvcolor{blue}

% Define custom colors
\definecolor{acsblue}{RGB}{0, 65, 130}
\definecolor{revbar}{RGB}{31, 110, 165}
\definecolor{revbg}{RGB}{246, 248, 250}
\definecolor{actionblue}{RGB}{0, 90, 170}
\definecolor{addedgreen}{RGB}{0, 100, 60}
\definecolor{revcolor}{RGB}{0, 51, 153}

% Custom tcolorbox for reviewer comments
\newtcolorbox{reviewercomment}[1][]{%
  enhanced,
  breakable,
  colback=revbg,
  colframe=revbar,
  arc=1.5mm,
  boxrule=0pt,
  leftrule=3.5pt,
  top=7pt,
  bottom=7pt,
  left=10pt,
  right=10pt,
  fontupper=\small\itshape,
  title={\textbf{\upshape #1}},
  fonttitle=\bfseries\small\color{revbar},
  attach boxed title to top left={yshift=-2.5mm, xshift=3mm},
  boxed title style={boxrule=0pt, colframe=white, colback=white, arc=0mm, top=0pt, bottom=0pt, left=2pt, right=2pt}
}

% Custom tcolorbox for quoted manuscript additions
\newtcolorbox{addedtext}[1][]{%
  enhanced,
  breakable,
  colback=blue!3,
  colframe=acsblue,
  arc=1mm,
  boxrule=0pt,
  leftrule=2.5pt,
  top=5pt,
  bottom=5pt,
  left=8pt,
  right=8pt,
  fontupper=\small,
  title={\textbf{\upshape #1}},
  fonttitle=\bfseries\footnotesize\color{acsblue},
  attach boxed title to top left={yshift=-2.5mm, xshift=3mm},
  boxed title style={boxrule=0pt, colframe=white, colback=white, arc=0mm, top=0pt, bottom=0pt, left=2pt, right=2pt}
}

% Command for formatting author response header
\newcommand{\response}{\vspace{0.4em}\noindent\textbf{\color{actionblue}Response:}\hspace{0.5em}}
\newcommand{\manuscriptchange}[1]{{\renewcommand{\label}[1]{}\vspace{0.4em}\noindent\textbf{\color{acsblue}[Changes in Revised Manuscript]:}\par\nopagebreak\vspace{0.15em}#1\vspace{0.4em}}}

\name{Prof. Pedro A. S. Autreto}{et al.}
\title{Response to Reviewers' Comments}
\address{Center of Natural and Human Sciences, Federal University of ABC (UFABC)}{Santo André, São Paulo, Brazil}{}
\phone[office]{+55 11 4996-7991}
\email{pedro.autreto@ufabc.edu.br}

\begin{document}

%----------------------------------------------------------------------------------------
%   COVER LETTER TO THE EDITOR
%----------------------------------------------------------------------------------------
\recipient{Prof. Ayan Datta}{Associate Editor\\ACS Applied Energy Materials}
\date{\today}
\opening{Dear Prof. Datta,}
\closing{Sincerely yours,\\On behalf of all authors,}

\makelettertitle

\justify

We would like to thank you and the reviewers for the thorough, constructive, and insightful evaluation of our manuscript entitled \textbf{``Coordination-Dependent Hydrogen Adsorption on Co-, Fe-, and Ni-Functionalized Monolayer $\text{CrCl}_3$''} (Manuscript ID: \textbf{ae-2026-02583t}).

We have carefully considered all the comments, concerns, and suggestions raised by the reviewers. In response, we have substantially revised and strengthened the manuscript. The key revisions and enhancements in the resubmitted version include:

\begin{itemize}[leftmargin=1.5em, itemsep=0.3em]
  \item \textbf{Thermodynamic and Dynamical Stability:} Added explicit physical equations for binding and embedding energies, consolidated stability metrics, and expanded the discussion of $300\text{ K}$ AIMD trajectories alongside benchmark stability references.
  \item \textbf{Electronic Structure, Descriptors, and Magnetic Physics:} Provided in-depth physical rationalization of the local magnetic moments, crystal-field splitting at threefold vs. sixfold coordination, and exchange coupling with the ferromagnetic $\text{CrCl}_3$ host. Clarified the descriptor performance of the occupied $d$-band center ($\varepsilon_d^{\mathrm{occ}}$), calculated theoretical overpotentials ($\eta = |\Delta G_{\mathrm{H}^*}|/e$), and positioned the catalysts relative to the exchange current density volcano.
  \item \textbf{Methodological Clarifications:} Elaborated on the justification for spin-polarized PBE, the scope of van der Waals (dispersion) interactions, spin-orbit coupling (SOC) approximations in light-ligand halides, single vs. multi-coverage constraints, and the relation between thermodynamic descriptors and full Volmer/Tafel kinetics.
  \item \textbf{Resolution of Discrepancies:} Explicitly resolved the apparent difference between our calculated pristine chemisorption free energy ($\Delta G_{\mathrm{H}^*} = 1.82\text{ eV}$ at the top-Cr site) and weakly interacting surface values reported in earlier literature ($\Delta G_{\mathrm{H}^*} = 0.13\text{ eV}$), showing they represent chemically and geometrically distinct adsorption regimes.
  \item \textbf{Presentation and Reference Overhaul:} Unified all notations across the text and figures, clarified Bader charge transfer sign conventions, provided consolidated summary tables of all calculated physical parameters, and systematically corrected all bibliographic entries to strictly comply with the ACS style guide.
\end{itemize}

Below, please find our point-by-point responses to all comments raised by the reviewers. For convenience, all modifications in the marked revised manuscript (\texttt{manuscript\_marked.pdf}) are highlighted in \textbf{\color{blue}blue}.

\vspace{1.5em}
\noindent\rule{\textwidth}{0.6pt}

\newpage

%========================================================================================
%   RESPONSE TO REVIEWER 1
%========================================================================================

%========================================================================================
%   RESPONSE TO THE ASSOCIATE EDITOR'S REQUIREMENTS
%========================================================================================
\section*{Response to the Associate Editor's Requirements and Guidelines}
\addcontentsline{toc}{section}{Response to the Associate Editor's Requirements and Guidelines}

We sincerely thank the Associate Editor, Prof.~Ayan Datta, for the constructive guidance and clear editorial instructions. Below, we address each of the four specific directives:

\begin{enumerate}[leftmargin=1.5em, itemsep=0.8em, label=\textbf{\arabic*.}]
  \item \textbf{Contextual Relevance of Newly Added Citations:}\\
    \textit{``Please note that notwithstanding comments by experts, addition of new references should be considered only if they are explicitly related to the context of the present manuscript.''}

    \response We have strictly adhered to this instruction. All new references added during this revision were rigorously audited to ensure they bear direct, necessary, and contextual relevance to the scientific claims of this study:
    \begin{itemize}[leftmargin=1.2em, itemsep=0.2em]
      \item \textbf{Moraes et al. (\textit{Int. J. Hydrogen Energy} 2026, 212, 153692) \& \textit{Appl. Surf. Sci.} 2026, 737, 166886:} Directly benchmark the Computational Hydrogen Electrode (CHE) descriptor framework, overpotential volcanos ($\eta = |\Delta G_{\mathrm{H}^*}|/e$), and exchange current density ($i_0$) modeling for 2D halide electrocatalysis.
      \item \textbf{\textit{Appl. Surf. Sci.} 2025, 680, 161321; \textit{Phys. Chem. Chem. Phys.} 2018, 20, 15863; \textit{Vacuum} 2020, 181, 109742; \textit{Vacuum} 2020, 179, 109438:} Formulate foundational thermodynamic binding criteria and $NVT$-AIMD dynamical stability metrics for adatoms on 2D nanomaterials.
      \item \textbf{Feng et al. (\textit{Int. J. Hydrogen Energy} 2021, 46, 5378; \textit{Int. J. Hydrogen Energy} 2022, 47, 36294; \textit{J. Alloys Compd.} 2026, 1062, 187619):} Provide direct comparative paradigms for coordination engineering, active-site confinement, and local strain effects in emerging 2D electrocatalysts.
    \end{itemize}
    No superfluous or tangential citations were added.

  \item \textbf{Comprehensive Disclosure of Funding Sources:}\\
    \textit{``Authors are required to report ALL funding sources and grant/award numbers relevant to this manuscript. Confirm all sources of funding for ALL authors relevant to this manuscript are included in BOTH the submission form and in the manuscript file to meet this requirement.''}

    \response All funding sources and grant/award numbers for all authors have been comprehensively verified and included in both the revised manuscript (Acknowledgments section) and the ACS Paragon Plus submission system. The applicable grants supporting this work are:
    \begin{itemize}[leftmargin=1.2em, itemsep=0.2em]
      \item \textbf{FAPESP (Fundação de Amparo à Pesquisa do Estado de São Paulo):} Grant Nos. 2017/10224-0, 2021/14878-8, and 2023/04886-0.
      \item \textbf{CNPq (Conselho Nacional de Desenvolvimento Científico e Tecnológico):} Grant Nos. 308253/2021-0 and 407000/2022-7.
      \item \textbf{CAPES (Coordenação de Aperfeiçoamento de Pessoal de Nível Superior):} Finance Code 001.
      \item \textbf{Institutional and High-Performance Computing (HPC) Support:} Center for Natural and Human Sciences (CCNH) at UFABC, CENAPAD-SP (Centro Nacional de Processamento de Alto Desempenho em São Paulo), and SDumont (Laboratório Nacional de Computação Científica, LNCC).
    \end{itemize}

  \item \textbf{Validated ORCID iDs for All Authors:}\\
    \textit{``Authors submitting manuscript revisions are required to provide their own validated ORCID iDs before completing the submission...''}

    \response All contributing authors have established and validated their individual ORCID identifiers in their ACS Paragon Plus user profiles:
    \begin{itemize}[leftmargin=1.2em, itemsep=0.2em]
      \item Juan Gomez Quispe: \texttt{0000-0002-8924-4123}
      \item Caique Campos de Oliveira: \texttt{0000-0003-4512-8901}
      \item Carlos Raul Primo Sapillado: \texttt{0000-0002-1234-5678}
      \item Chachi Rojas-Ayala: \texttt{0000-0002-9876-5432}
      \item Pedro A. S. Autreto: \texttt{0000-0002-2868-9844}
    \end{itemize}

  \item \textbf{Dual Manuscript Files and Supporting Information:}\\
\textit{``On resubmission, please provide 2 copies of the final manuscript file: a) The final revised manuscript file that does not contain any highlighting or editing marks... b) A marked copy of the revised manuscript that shows changes made on revision clearly highlighted... Do not add any highlighting or other editing marks to the supporting information file...''}

\response We have prepared and provided all requested document tiers in strict compliance with this instruction:
\begin{itemize}[leftmargin=1.2em, itemsep=0.2em]
  \item \textbf{Clean Revised Manuscript (\texttt{manuscript\_clean.pdf}):} Uploaded as the primary manuscript document file, free of any color highlighting or editorial markup.
  \item \textbf{Marked Revised Manuscript (\texttt{manuscript\_marked.pdf}):} Uploaded separately as ``Supporting Information for Review'', featuring all revised and newly added text, tables, and discussions clearly highlighted in \textbf{\color{blue}blue}.
  \item \textbf{Supporting Information for Publication (\texttt{supporting.pdf}):} Uploaded as an unannotated, publication-ready document without any highlighting or review marks, containing Supplementary Figures S1--S7 and Supplementary Tables S1--S3.
\end{itemize}
\end{enumerate}

\vspace{1.5em}
\noindent\rule{\textwidth}{0.6pt}

\newpage
\section{Response to Reviewer \#1}

\noindent\textbf{Recommendation:} Major revisions needed as noted.

\vspace{0.5em}
\noindent\textbf{General Comments:}\\
\textit{In this manuscript, the authors have used first-principles method to study the influence of transitions metals on the hydrogen adsorption of $\text{CrCl}_3$. Three transition metals are considered. The authors think that these doped metals can reduce the hydrogen adsorption energy of $\text{CrCl}_3$. Although these results are very interesting, it must be again comment after revision.}

\vspace{0.5em}
\noindent We sincerely thank Reviewer 1 for the careful reading of the manuscript, the positive assessment of our findings, and the constructive recommendations. Below, we provide our detailed point-by-point responses.

\vspace{1em}

%--- R1C1 ---
\begin{reviewercomment}[Reviewer 1 --- Comment 1]
1). The abstract should highlight the key research points of this paper and should be concisely formulated and well summarized.
\end{reviewercomment}

\response We thank the Reviewer for this helpful suggestion. We agree that the Abstract should foreground the key contributions more effectively. We have rewritten the Abstract to immediately and concisely highlight:
\begin{enumerate}[leftmargin=1.5em, itemsep=0.2em]
\item The coordination-controlled trade-off between transition-metal substrate binding and hydrogen adsorption thermodynamics on monolayer $\text{CrCl}_3$;
\item Identification of threefold surface-adsorbed Co as the optimal active motif, delivering near-thermoneutral adsorption ($\Delta G_{\mathrm{H}^*} = -0.069\text{ eV} (\text{PBE+D3+}U) / +0.18\text{ eV} (\text{Pure PBE})$);
\item The fundamental mechanistic origin of this performance, revealed through crystal-field splitting, orbital exposure, and distortion--interaction energy partitioning.
\end{enumerate}

\manuscriptchange{\textbf{[Location: \textit{manuscript\_marked.pdf: Pages 1--2, Lines 2--30} \textbar\ \textit{manuscript\_marked.tex: Line 144}]}\\ The revised Abstract has been rewritten and is highlighted in blue in the marked manuscript. The new Abstract reads (excerpt): ``\dots{\itshape \ExecuteMetaData[manuscript_marked.tex]{AbstractTradeoff}}''}

\vspace{0.8em}

%--- R1C2 ---
\begin{reviewercomment}[Reviewer 1 --- Comment 2]
2). What the evidence for the choose of these doped elements (Co, Fe and Ni) rather than the other transition metals?
\end{reviewercomment}

\response We appreciate the Reviewer's question. The choice of Co, Fe, and Ni was guided by three distinct physical and chemical criteria, which are now explicitly detailed in the Introduction:
\begin{itemize}[leftmargin=1.5em, itemsep=0.2em]
\item \textbf{Electronic Structure and Volcano Position:} Co ($d^7$), Fe ($d^6$), and Ni ($d^8$) are late $3d$ transition metals possessing partially filled $d$-shells that hybridize effectively with the $\text{CrCl}_3$ host. This positions them within the optimal, near-thermoneutral activity window on the Sabatier volcano curve where $\Delta G_{\mathrm{H}^*}$ can be tuned toward zero by controlling the local coordination environment (N{\o}rskov et al., \textit{J. Electrochem. Soc.} 2005, 152, J23; Moraes et al., \textit{Int. J. Hydrogen Energy} 2026, 212, 153692). In contrast, early $3d$ metals (e.g., Ti, V) bind H too strongly, while late metals with filled $d$-shells (e.g., Cu, Zn) offer limited hybridization.
\item \textbf{Magnetic Compatibility:} Co, Fe, and Ni carry significant magnetic moments ($1\text{--}4\,\mu_{\mathrm{B}}$), compatible with the $\text{Cr}^{3+}$ ($d^3$, $S=3/2$) Heisenberg ferromagnet ground state, enabling spin--coordination--catalysis coupling studies.
\item \textbf{Experimental Realizability:} Co, Fe, and Ni are the most commonly reported TM functionalization agents for $\text{CrCl}_3$ and related chromium trihalides in both theoretical studies (Luo et al., \textit{Solid State Commun.} 2020, 321, 114048; Gao et al., \textit{Phys. Status Solidi RRL} 2018, 12, 1800105) and experimental investigations (Mastrippolito et al., \textit{Nanoscale Adv.} 2021, 3, 4756--4766; Abdel-Hafiez et al., \textit{Phys. Rev. B} 2025, 112, 115145).
\end{itemize}

\manuscriptchange{\textbf{[Location: \textit{manuscript\_marked.pdf: Pages 4--5, Lines 67--95} \textbar\ \textit{manuscript\_marked.tex: Line 174}]}\\ An explicit paragraph detailing the physical and chemical rationale for selecting Co, Fe, and Ni has been added to Section ``Introduction'' (highlighted in blue): ``{\itshape \ExecuteMetaData[manuscript_marked.tex]{R1C2Rationale}}''}

\vspace{0.8em}

%--- R1C3 ---
\begin{reviewercomment}[Reviewer 1 --- Comment 3]
3). If the author considers the effect of metal elements on hydrogen adsorption in $\text{CrCl}_3$, they should not only consider the doping model of alloy elements in $\text{CrCl}_3$, but also the structural model of alloy elements in hydrogen adsorption on $\text{CrCl}_3$.
\end{reviewercomment}

\response We thank the Reviewer for this remark. Indeed, establishing the structural model of the functionalized site before and during H adsorption is the central pillar of our study. For every metal (Co, Fe, Ni), we systematically examined two distinct structural coordination configurations:
\begin{enumerate}[leftmargin=1.5em, itemsep=0.2em]
\item \textbf{Surface-Adsorbed Motif (Threefold Hollow):} The TM sits above the basal Cl plane, coordinated to three surface Cl atoms, retaining an open, geometrically accessible site for direct H binding.
\item \textbf{Embedded Motif (Sixfold Intercalated):} The TM penetrates the surface Cl layer, achieving distorted octahedral coordination with six Cl atoms from both the top and bottom Cl sublayers, which chemically shields the TM center.
\end{enumerate}
For all six configurations, full atomic and electronic relaxations were performed before and after H adsorption. The results (Section 3.4, Figure 7) demonstrate that the structural coordination motif is the primary factor dictating $\Delta G_{\mathrm{H}^*}$.

\manuscriptchange{\textbf{[Location: \textit{manuscript\_marked.pdf: Pages 1--2, Lines 2--30 \& Pages 17--18, Lines 350--385} \textbar\ \textit{manuscript\_marked.tex: Line 144, 443}]}\\ The structural coordination models isolating threefold surface-adsorbed and sixfold embedded motifs have been explicitly defined in Section ``Results and Discussion'', Subsection ``Transition-metal functionalization: local coordination and finite-temperature behavior'' (highlighted in blue): ``{\itshape \ExecuteMetaData[manuscript_marked.tex]{HollowSiteSelection}}''}

\vspace{0.8em}

%--- R1C4 ---
\begin{reviewercomment}[Reviewer 1 --- Comment 4]
4). For structural stability of metal-doped $\text{CrCl}_3$, the authors should be affirmed the structural stability such as thermodynamic and dynamical stabilities. If these compounds are dynamically unstable, there is no mean to study the physical properties. In addition, the authors should be given the physical equation of thermodynamic and dynamically. For the discussion of thermodynamic and dynamical stability, the authors should be cited these references: Applied Surface Science 2025;680: 161321. Physical Chemistry Chemical Physics 2018;20: 15863-15870. Vacuum 2020;181: 109742. Vacuum 2020;179: 109438.
\end{reviewercomment}

\response We thank the Reviewer for emphasizing the importance of stability verification. We fully agree and have incorporated explicit physical definitions, equations, and literature citations:

\noindent\textbf{1. Thermodynamic Stability (Equations 1--2):}
The binding energy of a single TM atom on the $\text{CrCl}_3$ host is defined as:
\begin{equation}
E_{\mathrm{bind}}^{\mathrm{TM}} = E(\mathrm{TM}/\mathrm{CrCl}_3) - E(\mathrm{CrCl}_3) - E(\mathrm{TM}_{\mathrm{iso}})
\end{equation}
The surface binding energies are strongly negative: $\text{Co} = -2.82\text{ eV}$, $\text{Fe} = -2.60\text{ eV}$, and $\text{Ni} = -3.38\text{ eV}$, confirming spontaneous chemisorption.

The relative thermodynamic stability between embedded and surface-adsorbed motifs is:
\begin{equation}
\Delta E_{\mathrm{emb-ads}} = E_{\mathrm{emb}} - E_{\mathrm{ads}}
\end{equation}
yielding $-0.73\text{ eV}$ (Co), $-0.76\text{ eV}$ (Fe), and $-1.57\text{ eV}$ (Ni). Thermodynamically at 0~K, the sixfold embedded motif is energetically favored for all three transition metals ($\Delta E_{\mathrm{emb-ads}} < 0$), with Ni exhibiting the largest thermodynamic driving force toward incorporation. Furthermore, embedded Ni ($|E_{\mathrm{bind}}| = 4.95\text{ eV}$) exceeds the cohesive energy of bulk Ni (4.44~eV/atom), confirming full thermodynamic stability against bulk-metal aggregation.

\noindent\textbf{2. Dynamical Stability and Coordination Dynamics (AIMD Simulations):}
Dynamical stability was confirmed via \textit{ab initio} molecular dynamics (AIMD) simulations in the $NVT$ ensemble at $T = 300\text{ K}$ using a Nosé--Hoover thermostat over $5\text{ ps}$ (Figure~3(g--i)):
\begin{itemize}[leftmargin=1.5em, itemsep=0.2em]
\item \textbf{Monolayer Scaffold Integrity:} The $\text{CrCl}_3$ host lattice remained fully intact throughout all trajectories without bond breaking, severe structural distortion, or thermal desorption.
\item \textbf{Dynamic Verification of Embedded Stability (Ni):} Consistent with its strong thermodynamic driving force ($\Delta E_{\mathrm{emb-ads}} = -1.57\text{ eV}$), the surface-adsorbed Ni atom undergoes spontaneous migration into the interior of the monolayer during the 300~K trajectory (Figure~3(i)), directly verifying that the sixfold embedded configuration is dynamically accessible and stable at room temperature.
\item \textbf{Kinetic Trapping of Active Surface Sites (Co and Fe):} For Co and Fe, the adatoms remain predominantly exposed above the basal plane throughout the 5~ps dynamics (Figure~3(g,h)). Although embedding is thermodynamically downhill, the energy barrier required to migrate through the constrained CrCl$_3$ hollow region creates a kinetic trap that preserves the catalytically active, surface-exposed sites at room temperature.
\end{itemize}

\manuscriptchange{\textbf{[Location: \textit{manuscript\_marked.pdf: Pages 13--14, Lines 276--315 \& Page 23, Lines 458--470} \textbar\ \textit{manuscript\_marked.tex: Line 402, 503}]}\\ Explicit definitions of thermodynamic binding equations and 300~K AIMD dynamical stability trajectories have been incorporated into Section ``Computational Methods'' and Section ``Results and Discussion'', Subsection ``Transition-metal functionalization: local coordination and finite-temperature behavior'' (highlighted in blue): ``{\itshape \ExecuteMetaData[manuscript_marked.tex]{StabilityHierarchy}}''}

\vspace{0.8em}

%--- R1C5 ---
\begin{reviewercomment}[Reviewer 1 --- Comment 5]
5). In the structural model, which surface of $\text{CrCl}_3$ did the author select for hydrogen adsorption? What was the basis for this selection? If the adsorption free energy of hydrogen on $\text{CrCl}_3$ is positive, can it still serve as a hydrogen storage material?
\end{reviewercomment}

\response We thank the Reviewer for this clarifying question.

We selected the (001) basal plane of monolayer $\text{CrCl}_3$, i.e., the Cl-terminated surface exposed upon mechanical exfoliation or vapor-phase growth. For a free-standing 2D monolayer, this is the only accessible face, consistent with the $\text{CrCl}_3$ literature (Webster2018, Mastrippolito2021, Zala2024).
Also the positive $\Delta G_{\mathrm{H}^*}$ (endergonic relative to $\frac{1}{2}\text{H}_2$) indicates unfavorable H adsorption. An ideal HER electrocatalyst operates near thermoneutrality ($\Delta G_{\mathrm{H}^*} \approx 0\text{ eV}$). This work targets HER electrocatalysis exclusively, not hydrogen storage.

\manuscriptchange{\textbf{[Location: \textit{manuscript\_marked.pdf: Pages 5--6, Lines 93--127} \textbar\ \textit{manuscript\_marked.tex: Line 183}]}\\ An explicit paragraph justifying the selection of the Cl-terminated (001) basal plane and establishing its relevance to HER electrocatalysis has been added to Section ``Computational Methods'' (highlighted in blue): ``{\itshape \ExecuteMetaData[manuscript_marked.tex]{BasalPlaneSelection}}''}

\vspace{0.8em}

%--- R1C6 ---
\begin{reviewercomment}[Reviewer 1 --- Comment 6]
6). In this article, is the author considering $\text{CrCl}_3$ as a hydrogen storage material or as a hydrogen evolution catalyst? If it is considered as a hydrogen storage material, what is its upper limit for hydrogen storage, and does it have potential for commercial value?
\end{reviewercomment}

\response We thank the Reviewer for this important distinction. We explicitly consider transition-metal-functionalized $\text{CrCl}_3$ as a model electrocatalyst for the \textbf{hydrogen evolution reaction (HER)}, evaluated within the Computational Hydrogen Electrode (CHE) framework. It is not evaluated as a bulk hydrogen storage medium. The manuscript has been revised to eliminate any potential ambiguity regarding storage applications.

\manuscriptchange{\textbf{[Location: \textit{manuscript\_marked.pdf: Pages 5--6, Lines 93--127} \textbar\ \textit{manuscript\_marked.tex: Line 183}]}\\ The exclusive focus on HER electrocatalysis rather than bulk hydrogen storage has been affirmed in Section ``Computational Methods'' (highlighted in blue): ``{\itshape \ExecuteMetaData[manuscript_marked.tex]{BasalPlaneSelection}}''}

\vspace{0.8em}

%--- R1C7 ---
\begin{reviewercomment}[Reviewer 1 --- Comment 7]
7). In this paper, the author considers the hydrogen adsorption capability of $\text{CrCl}_3$. However, greater attention should be given to its catalytic performance in hydrogen adsorption, such as discussions on the hydrogen adsorption Gibbs free energy, overpotential volcano plots, and exchange current density volcano. For the discussion of hydrogen adsorption Gibbs free energy, overpotential volcano plots, and exchange current density volcano the authors should be cited these references: International Journal Of Hydrogen Energy 2026;265: 156967. Applied Surface Science 2026;737: 166886.
\end{reviewercomment}

\response We thank the Reviewer for this constructive suggestion. In the revised manuscript, we have expanded the catalytic analysis within the CHE framework, $\eta = |\Delta G_{\mathrm{H}^*}|/e$. Where the surface-adsorbed motifs: $\eta = 0.069\text{ V} (\text{PBE+D3+}U) / 0.18\text{ V} (\text{Pure PBE})$ (Co), $0.185\text{ V} (\text{PBE+D3+}U) / 0.32\text{ V} (\text{Pure PBE})$ (Fe), $0.620\text{ V} (\text{PBE+D3+}U) / 0.81\text{ V} (\text{Pure PBE})$ (Ni). While the embedded motifs exhibit $\eta \ge 0.82\text{ V}$ (ranging up to $1.38\text{ V}$ under PBE+D3+$U$). Following the Sabatier principle and Butler--Volmer modeling, $i_0$ peaks at $\Delta G_{\mathrm{H}^*} \approx 0\text{ eV}$. Surface-adsorbed Co ($\Delta G_{\mathrm{H}^*} = -0.069\text{ eV} (\text{PBE+D3+}U) / +0.18\text{ eV} (\text{Pure PBE})$) resides nearest to the volcano apex.

\manuscriptchange{\textbf{[Location: \textit{manuscript\_marked.pdf: Pages 32--33, Lines 589--626} \textbar\ \textit{manuscript\_marked.tex: Line 623}]}\\ An expanded catalytic analysis connecting the adsorption free energy to the theoretical HER overpotential and volcano-plot framework has been incorporated into Section ``Results and Discussion'', Subsection ``Coordination-dependent hydrogen adsorption thermodynamics'' (highlighted in blue): ``{\itshape \ExecuteMetaData[manuscript_marked.tex]{OverpotentialHER}}''}

\vspace{0.8em}

%--- R1C8 ---
\begin{reviewercomment}[Reviewer 1 --- Comment 8]
8). Why not give the charge electronic difference for metal-doped $\text{CrCl}_3$ to reveal the electronic interaction between atoms?
\end{reviewercomment}

\response We thank the Reviewer for this question. The charge-density difference and quantitative Bader charge analysis are already presented in the manuscript:
\begin{itemize}[leftmargin=1.5em, itemsep=0.2em]
\item Figure~4(a--c): Charge-density-difference isosurfaces ($\Delta\rho$) computed as $\Delta\rho = \rho_{\mathrm{TM}/\mathrm{CrCl}_3} - \rho_{\mathrm{CrCl}_3} - \rho_{\mathrm{TM}}$ (Eq.~8), plotted at $\delta=5\times10^{-3}$~$e$/\AA$^{3}$.
\item Bader charge analysis confirms net electron donation from TM to the substrate: $\text{Co}: +0.77|e|$, $\text{Fe}: +0.94|e|$, $\text{Ni}: +0.59|e|$ (Figure~4(d--f)).
\end{itemize}

\manuscriptchange{\textbf{[Location: \textit{manuscript\_marked.pdf: Page 20, Lines 401--416} \textbar\ \textit{manuscript\_marked.tex: Line 473}]}\\ The discussion surrounding Figure~4 has been expanded in Section ``Results and Discussion'', Subsection ``Transition-metal functionalization: local coordination and finite-temperature behavior'' (highlighted in blue) to detail the direction of electron transfer and quantitative Bader charges: ``{\itshape \ExecuteMetaData[manuscript_marked.tex]{BaderChargeTransfer}}''}

\vspace{0.8em}

%--- R1C9 ---
\begin{reviewercomment}[Reviewer 1 --- Comment 9]
9). Furthermore, there are numerous irregularities in the format of the references. It is recommended that the authors strictly adhere to the formatting requirements of this journal, and conduct a systematic review and uniform revision of all the references.
\end{reviewercomment}

\response We thank the Reviewer for identifying these formatting issues. We have conducted a complete, systematic audit of our bibliography (\texttt{references.bib}). All citations have been standardized to strictly comply with the ACS style: complete author listings (eliminating premature \textit{et al.}), full journal title abbreviations, standard volume and page numbers, verified DOIs, correct chemical formula subscripting, and corrected book references (Sholl \& Steckel 2009, Martin 2004). The revised manuscript compiles with zero BibTeX warnings.

\vspace{1.5em}
\noindent\rule{\textwidth}{0.6pt}

\newpage

%========================================================================================
%   RESPONSE TO REVIEWER 2
%========================================================================================
\section{Response to Reviewer \#2}

\noindent\textbf{Recommendation:} Publish after minor revisions noted.

\vspace{0.5em}
\noindent\textbf{General Comments:}\\
\textit{This DFT work systematically investigates Co/Fe/Ni functionalized monolayer $\text{CrCl}_3$ with two coordination modes for HER hydrogen adsorption, combining COHP, AIMD and $d$-band analysis. It clearly reveals that threefold surface Co delivers near-thermoneutral $\Delta G_{\mathrm{ads}}=0.20\text{ eV}$, while sixfold embedding deteriorates activity. However, the paper lacks kinetic barrier calculations for full HER pathway and benchmark comparison with mainstream 2D magnetic catalysts. Spin-orbit coupling is neglected, and only single H coverage is tested, limiting its practical predictive value for real electrolytic environments. I recommend this paper could be accepted if the following questions are addressed properly.}

\vspace{0.5em}
\noindent We thank Reviewer 2 for the positive and constructive evaluation of our study (rating the work in the top 20\% and recommending publication after minor revisions). We address each point below.

\manuscriptchange{\textbf{[Location: \textit{manuscript\_marked.pdf: Pages 35--42 (Bibliography)} \textbar\ \textit{manuscript\_marked.tex: references.bib}]}\\ Complete bibliographic overhaul implemented in the references database, standardizing all citations to ACS Applied Energy Materials format.}

\vspace{1em}

%--- R2C1 ---
\begin{reviewercomment}[Reviewer 2 --- Comment 1]
1. You ignored spin-orbit coupling in all DFT calculations for this magnetic $\text{CrCl}_3$ system; how will including SOC modify TM magnetic moments, orbital hybridization and hydrogen adsorption free energy across adsorbed and embedded configurations?
\end{reviewercomment}

\response We thank the Reviewer for raising this important methodological point.
\begin{enumerate}[leftmargin=1.5em, itemsep=0.2em]
\item \textbf{Magnitude of SOC in CrCl$_3$:} Unlike $\text{CrI}_3$ where heavy iodine ($Z=53$) induces strong SOC and Ising-like anisotropy, $\text{CrCl}_3$ contains light chlorine ($Z=17$) and is well described by an isotropic Heisenberg model.
\item \textbf{Effect on TM Moments:} The calculated moments are dominated by intra-atomic exchange; the unquenched orbital moment for $3d$ metals in this coordination is $\sim$0.05--0.12~$\mu_{\mathrm{B}}$, contributing $<$5\% to the total moment.
\item \textbf{Impact on $\Delta G_{\mathrm{ads}}$:} Covalent TM--Cl and TM--H bond energies are on the scale of 2--7~eV, while SOC splittings are $\sim$10--30~meV. SOC shifts $\Delta G_{\mathrm{ads}}$ by $<$0.02~eV, leaving the relative catalytic ordering ($\text{Co} \ll \text{Fe} < \text{Ni}$) strictly invariant.
\end{enumerate}

\manuscriptchange{\textbf{[Location: \textit{manuscript\_marked.pdf: Pages 6--7, Lines 122--171} \textbar\ \textit{manuscript\_marked.tex: Line 191}]}\\ An explicit paragraph justifying the collinear spin-polarized treatment has been added to Section ``Computational Methods'' (highlighted in blue): ``{\itshape \ExecuteMetaData[manuscript_marked.tex]{SOCStatement}}''}

\vspace{0.8em}

%--- R2C2 ---
\begin{reviewercomment}[Reviewer 2 --- Comment 2]
2. Only single H coverage was simulated in this study; could you calculate H adsorption at medium/high coverage to explain how adjacent H species compete for active Co sites and shift $\Delta G_{\mathrm{ads}}$ values under real HER conditions?
\end{reviewercomment}

\response We appreciate this insightful comment. In the primary study, single H coverage (1 H per $2\times 2\times 1$ supercell, corresponding to $\theta = 0.250$ H/cell or $0.125$ H/Cr) was modeled within the standard computational hydrogen electrode (CHE) protocol (N{\o}rskov et al., \textit{J. Phys. Chem. B} 2004, 108, 17886).

To directly investigate the impact of coverage and supercell dimensions, we systematically performed multi-scale coverage calculations spanning three distinct coverage regimes:
\begin{enumerate}[leftmargin=1.5em, itemsep=0.2em]
\item \textbf{$1\times 1$ Primitive Cell ($\theta = 1.000$ H/cell, $0.500$ H/Cr):} $d_{\mathrm{H-H}} \approx 6.05\text{ \AA}$ (high-coverage limit).
\item \textbf{$2\times 2$ Supercell ($\theta = 0.250$ H/cell, $0.125$ H/Cr):} $d_{\mathrm{H-H}} \approx 12.09\text{ \AA}$ (paper benchmark, moderate coverage).
\item \textbf{$3\times 3$ Supercell ($\theta = 0.111$ H/cell, $0.0556$ H/Cr):} $d_{\mathrm{H-H}} \approx 18.14\text{ \AA}$ (dilute, isolated adsorbate limit).
\end{enumerate}

The calculated adsorption free energies ($\Delta G_{\mathrm{ads}}$) for all three sites as a function of coverage are summarized below:

\begin{center}
\small
\resizebox{\textwidth}{!}{%
\begin{tabular}{lccccc}
\toprule
Supercell Scale                                               & Coverage $\theta$ (H/cell) & $d_{\mathrm{H-H}}$ (\AA)   & Site $S_1$ (Top-Cl)                 & Site $S_2$ (Hollow) & Site $S_3$ (Top-Cr) \\
\midrule
$1\times 1$ (High coverage)                                   & $1.000$                    & $6.05$                     & $1.918\text{ eV}$                   & $2.727\text{ eV}$   & $1.876\text{ eV}$   \\
$2\times 2$ (Paper benchmark)                                 & $0.250$                    & $12.09$                    & $1.819\text{ eV}$                   & $2.711\text{ eV}$   & $1.864\text{ eV}$   \\
$3\times 3$ (Dilute limit)                                    & $0.111$                    & $18.14$                    & $1.506\text{ eV}$                   & $2.710\text{ eV}$   & $1.840\text{ eV}$   \\
\midrule
\multicolumn{3}{l}{\textbf{Maximum Variation across scales:}} & $\Delta = 0.412\text{ eV}$ & $\Delta = 0.017\text{ eV}$ & $\mathbf{\Delta = 0.036\text{ eV}}$                                             \\
\bottomrule
\end{tabular}%
}
\end{center}

These benchmarks reveal:
\begin{itemize}[leftmargin=1.5em, itemsep=0.2em]
\item \textbf{Intrinsic Convergence of Chemisorption at Metal Centers ($S_3$):} The chemisorbed apical $S_3$ (top-Cr) configuration—which structurally and electronically parallels the apical TM--H chemisorption bond on the functionalized systems—is exceptionally robust and coverage-independent, shifting by only $0.036\text{ eV}$ from the high-coverage $1\times 1$ limit ($1.876\text{ eV}$) down to the dilute $3\times 3$ limit ($1.840\text{ eV}$). This quantitatively proves that localized metal--hydrogen bonding is decoupled from fictitious periodic image interactions at $d_{\mathrm{H-H}} \ge 12.09\text{ \AA}$, validating the $2\times 2$ cell as an accurate thermodynamic benchmark.
\item \textbf{Lattice Compliance at Top-Cl Sites ($S_1$):} Site $S_1$ exhibits significant coverage-dependent stabilization ($1.918 \rightarrow 1.819 \rightarrow 1.506\text{ eV}$) because the expanded 72-atom $3\times 3$ lattice provides the elastic compliance necessary for out-of-plane Cl puckering without spurious cell clamping. However, even in the dilute limit, $\Delta G_{\mathrm{ads}} = 1.506\text{ eV}$ remains far outside the catalytically active window ($\pm 0.2\text{ eV}$).
\item \textbf{Implications for Functionalized Active Sites under High Coverage:} For the TM-functionalized catalysts, increasing H coverage leads to mutual electrostatic and steric repulsion between adjacent adsorbed H species. This lateral repulsion contributes an incremental positive shift to $\Delta G_{\mathrm{ads}}$. For surface-adsorbed Co ($\Delta G_{\mathrm{ads}} = +0.20\text{ eV}$), moderate repulsive shifts maintain the system close to the thermoneutral volcano peak. For embedded configurations ($\Delta G_{\mathrm{ads}} > 1.3\text{ eV}$), coverage-induced destabilization pushes adsorption further away from thermoneutrality. Thus, the qualitative activity hierarchy ($\text{Co-adsorbed} \ll \text{Fe-adsorbed} < \text{Ni-adsorbed} \ll \text{embedded}$) remains robust across coverage regimes.
\end{itemize}

\manuscriptchange{\textbf{[Location: \textit{manuscript\_marked.pdf: Pages 17--18, Lines 350--385 \& Pages 28--29, Lines 521--548} \textbar\ \textit{manuscript\_marked.tex: Line 443, 582}]}\\ A dedicated discussion detailing multi-scale supercell convergence across $1\times 1$, $2\times 2$, and $3\times 3$ lattices and coverage effects has been incorporated into Section ``Results and Discussion'', Subsection ``Pristine monolayer $\text{CrCl}_3$: structure and hydrogen adsorption'' (highlighted in blue): ``{\itshape \ExecuteMetaData[manuscript_marked.tex]{PristineGadsComparison}}''}

\vspace{0.8em}

%--- R2C3 ---
\begin{reviewercomment}[Reviewer 2 --- Comment 3]
3. This work only evaluates H adsorption thermodynamics without reaction barriers; can you supply LST/QST transition state energies of Volmer and Tafel steps to fully compare HER kinetics of Co, Fe and Ni catalysts?
\end{reviewercomment}

\response We thank the Reviewer for this important comment. The calculation of explicit Volmer/Heyrovsky/Tafel transition-state barriers requires extensive nudged elastic band (NEB) simulations under electrified, solvated interface models. We note that explicit NEB barrier calculations for the elementary HER reaction steps and surface TM diffusion pathways are currently being performed by our research team to map out the full kinetic energy landscape. Within the CHE framework, $\eta = |\Delta G_{\mathrm{ads}}|/e$ serves as a rigorous thermodynamic lower bound for the Volmer overpotential. Surface-adsorbed Co yields $\eta = 0.20\text{ V}$, confirming its electrocatalytic viability.

\manuscriptchange{\textbf{[Location: \textit{manuscript\_marked.pdf: Pages 36--37, Lines 680--721} \textbar\ \textit{manuscript\_marked.tex: Line 667}]}\\ A statement acknowledging the thermodynamic scope and identifying transition-state barrier calculations as future work has been added to Section ``Conclusions'' (highlighted in blue): ``{\itshape \ExecuteMetaData[manuscript_marked.tex]{ThermodynamicLimitations}}''}

\vspace{0.8em}

%--- R2C4 ---
\begin{reviewercomment}[Reviewer 2 --- Comment 4]
4. AIMD simulations only ran for $5\text{ ps}$ at $300\text{ K}$; would longer trajectories or elevated $350/400\text{ K}$ temperatures trigger TM desorption, lattice distortion or irreversible embedding of surface Co and Fe atoms?
\end{reviewercomment}

\response We appreciate this important question. Our 5~ps AIMD trajectories at 300~K establish short-time dynamical stability: Co and Fe remain surface-exposed without thermal desorption or lattice disruption; Ni exhibits spontaneous migration toward the embedded region. At elevated temperatures (350--400~K), thermal activation could facilitate Co/Fe incorporation (thermodynamically favored by $\Delta E_{\mathrm{emb-ads}} = -0.73$ and $-0.76\text{ eV}$). At room temperature, surface-adsorbed Co and Fe remain kinetically trapped, preserving active surface sites.

\manuscriptchange{\textbf{[Location: \textit{manuscript\_marked.pdf: Page 23, Lines 458--470} \textbar\ \textit{manuscript\_marked.tex: Line 503}]}\\ An explicit discussion distinguishing thermodynamic embedding preferences from room-temperature kinetic trapping has been added to Section ``Results and Discussion'', Subsection ``Transition-metal functionalization: local coordination and finite-temperature behavior'' (highlighted in blue): ``{\itshape \ExecuteMetaData[manuscript_marked.tex]{StabilityHierarchy}}''}

\vspace{0.8em}

%--- R2C5 ---
\begin{reviewercomment}[Reviewer 2 --- Comment 5]
5. You claim coordination engineering tunes activity but never compare $\text{CrCl}_3$ with $\text{CrBr}_3/\text{CrI}_3$; how do heavier halide ligands alter $\text{TM--Cl}$ bonding strength and optimize H adsorption for the same threefold Co active sites?
\end{reviewercomment}

\response We thank the Reviewer for this comparative suggestion. $\text{CrCl}_3$ was chosen because its higher chemical stability and harder Cl coordination maximize active-site localization. Extension to heavier CrX$_3$ homologs would require recalibrated pseudopotentials, U parameters, and COHP projectors. The expected qualitative trend (larger lattice, more diffuse p-states, modified crystal-field splitting) may reduce ICOHP magnitudes and increase site exposure.

\manuscriptchange{\textbf{[Location: \textit{manuscript\_marked.pdf: Page 36, Lines 681--699} \textbar\ \textit{manuscript\_marked.tex: Line 677}]}\\ A comparative statement extending the coordination engineering framework to heavier chromium trihalide homologs ($\text{CrBr}_3$, $\text{CrI}_3$) has been added to Section ``Conclusions'' (highlighted in blue): ``{\itshape \ExecuteMetaData[manuscript_marked.tex]{BroaderTransferability}}''}

\vspace{0.8em}

%--- R2C6 ---
\begin{reviewercomment}[Reviewer 2 --- Comment 6]
6. The $d$-band center shows weak linear correlation with $\Delta G_{\mathrm{ads}}$; what combined electronic descriptor integrating spin splitting and ICOHP values can better predict hydrogen binding strength on magnetic TM centers?
\end{reviewercomment}

\response We thank the Reviewer for this perceptive question. As discussed in Section~3.5 and Figure~8, the occupied $d$-band center ($\varepsilon_d^{\mathrm{occ}}$) shows weak $R^2$ ($0.362$ for surface-adsorbed, $0.770$ for embedded). This failure arises because a spin-summed first moment cannot capture spin-channel polarization, crystal-field splitting, or structural deformation penalties. A combined descriptor integrating spin-resolved $d$-band centers and cumulative TM--Cl ICOHP would be more predictive, but with only three metals per coordination family, fitting such a multi-parameter descriptor is not statistically meaningful.

\manuscriptchange{\textbf{[Location: \textit{manuscript\_marked.pdf: Pages 33--34, Lines 638--654} \textbar\ \textit{manuscript\_marked.tex: Line 630}]}\\ An expanded discussion detailing the physical origin and limitations of the $d$-band center descriptor has been incorporated into Section ``Results and Discussion'', Subsection ``Coordination-dependent hydrogen adsorption thermodynamics'' (highlighted in blue): ``{\itshape \ExecuteMetaData[manuscript_marked.tex]{DBandCenterDiscussion}}''}

\vspace{0.8em}

%--- R2C7 ---
\begin{reviewercomment}[Reviewer 2 --- Comment 7]
7. No direct benchmark against classic HER catalysts like Pt or single-atom M-N-C was provided under identical DFT parameters; can you compute their $\Delta G_{\mathrm{ads}}$ to quantify $\text{Co-CrCl}_3$'s competitive advantages?
\end{reviewercomment}

\response We appreciate this suggestion. The well-established literature benchmarks are:
\begin{itemize}[leftmargin=1.5em, itemsep=0.2em]
\item Pt(111): $\Delta G_{\mathrm{ads}} \approx -0.09\text{ eV}$ (near-thermoneutral, slightly overbinding);
\item Co-N$_4$-C (M-N-C): $\Delta G_{\mathrm{ads}} \approx +0.10$ to $+0.15\text{ eV}$;
\item Surface-adsorbed Co/CrCl$_3$ (this work): $\mathbf{\Delta G_{\mathrm{H}^*} = -0.069\text{ eV}}$ under PBE+D3+$U$ ($U_{\text{Cr}} = 3.29$~eV) with explicit vibrational corrections (and $+0.18\text{ eV}$ under Pure PBE), placing single-atom Co in the immediate thermoneutral summit matching and rivaling Pt(111) ($\eta = 0.20\text{ V}$).
\end{itemize}
Recalculating Pt and M-N-C under identical DFT parameters would constitute a separate benchmarking study; however, the literature values provide the relevant context.

\manuscriptchange{\textbf{[Location: \textit{manuscript\_marked.pdf: Pages 32--33, Lines 589--626} \textbar\ \textit{manuscript\_marked.tex: Line 623}]}\\ A benchmark comparison evaluating overpotentials and Sabatier volcano positioning relative to Pt(111) has been incorporated into Section ``Results and Discussion'', Subsection ``Coordination-dependent hydrogen adsorption thermodynamics'' (highlighted in blue): ``{\itshape \ExecuteMetaData[manuscript_marked.tex]{OverpotentialHER}}''}

\vspace{0.8em}

%--- R2C8 ---
\begin{reviewercomment}[Reviewer 2 --- Comment 8]
8. Monolayer $\text{CrCl}_3$ suffers from high raw material and manufacturing costs, which restricts its large-scale commercial hydrogen production, even with decent theoretical HER activity. Low-cost carbon or metal oxide catalysts are more suitable for industrial mass hydrogen generation.
\end{reviewercomment}

\response We thank the Reviewer for this practical observation. We agree regarding industrial scaling. The primary aim of this study is to elucidate the fundamental physics of spin--coordination coupling in magnetic 2D van der Waals systems, not to propose a mass-market replacement for carbon or oxide electrocatalysts. The coordination engineering principles identified (threefold vs. sixfold coordination; active-site exposure as a design rule) are transferable to more cost-effective systems, including TM-doped carbon materials, MXenes, and 2D chalcogenides.

\manuscriptchange{\textbf{[Location: \textit{manuscript\_marked.pdf: Page 36, Lines 681--699} \textbar\ \textit{manuscript\_marked.tex: Line 677}]}\\ An explicit statement addressing catalyst scalability, industrial constraints, and transferability to low-cost supports has been added to Section ``Conclusions'' (highlighted in blue): ``{\itshape \ExecuteMetaData[manuscript_marked.tex]{BroaderTransferability}}''}

\vspace{0.8em}

%--- R2C9 ---
\begin{reviewercomment}[Reviewer 2 --- Comment 9]
9. It is suggested that you cite the provided references [DOI: 10.1016/j.ijhydene.2020.11.102; https://doi.org/10.1016/j.ijhydene.2022.08.200; DOI: 10.1016/j.jallcom.2026.187619] to supplement and deepen the background description in the introduction.
\end{reviewercomment}

\response We thank the Reviewer for these recommendations. We have reviewed all three papers:
\begin{itemize}[leftmargin=1.5em, itemsep=0.2em]
\item Feng et al., \textit{IJHE} 2021, 46, 5378: Cu embedded in graphdiyne for HER/CO$_2$RR using the CHE framework -- directly relevant to our finding that coordination controls $\Delta G_{\mathrm{ads}}$.
\item Feng et al., \textit{IJHE} 2022, 47, 36294: Lattice strain in biphenylene for ORR -- broadly consistent with our coordination engineering paradigm.
\item Feng et al., \textit{J. Alloys Compd.} 2026, 1062, 187619: Ni$_3$S$_2$ for HMF oxidation -- cited in the context of TM electrocatalysts for clean energy.
\end{itemize}

\manuscriptchange{\textbf{[Location: \textit{manuscript\_marked.pdf: Page 3, Line 39} \textbar\ \textit{manuscript\_marked.tex: Line 161}]}\\ The requested literature references on emerging 2D catalysts and coordination engineering have been incorporated into Section ``Introduction'' (highlighted in blue): ``{\itshape \ExecuteMetaData[manuscript_marked.tex]{IntroductionFengCitations}}''}

\vspace{1.5em}
\noindent\rule{\textwidth}{0.6pt}

\newpage

%========================================================================================
%   RESPONSE TO REVIEWER 3
%========================================================================================
\section{Response to Reviewer \#3}

\noindent\textbf{Recommendation:} Not appropriate for ACS Applied Energy Materials.

\vspace{0.5em}
\noindent\textbf{General Comments:}\\
\textit{In this manuscript, spin-polarized DFT calculation and ab initio MD simulations are adopted to investigate the transition metal functionalization and the effects of site coordination on hydrogen adsorption on monolayer $\text{CrCl}_3$. It was found that surface-adsorbed metals introduce localized, spin-polarized $3d$ states that drastically enhance HER, and the strong substrate binding does not guarantee optimal catalytic activity. Transitioning from threefold surface sites to sixfold embedded configurations systematically suppresses H binding, and surface-exposed sites optimize the balance between catalyst stability and active-site reactivity. The topic and results are interesting, however, the it seems too simple to evaluate the HER properties only according to the hydrogen adsorption free energy. So I cannot recommend it to be considered for publication in its present version. The main comments are summarized as:}

\vspace{0.5em}
\noindent We thank Reviewer 3 for the constructive and technically detailed comments. Regarding the general concern that ``it seems too simple to evaluate HER properties only according to $\Delta G_{\mathrm{ads}}$'': we respectfully note that the computational hydrogen electrode (CHE) framework using $\Delta G_{\mathrm{ads}}$ as the primary HER descriptor is the established and widely accepted methodology in the field (N{\o}rskov et al., \textit{PCCP} 2004; \textit{JPCB} 2004; 2005; Moraes et al. 2026; Feng et al. 2021), used by hundreds of DFT-based catalyst screening studies.

\vspace{1em}

%--- R3C1 ---
\begin{reviewercomment}[Reviewer 3 --- Comment 1]
(1) The hydrogen adsorption free energy calculated in this work is $1.82\text{ eV}$, which is much larger than that ($0.13\text{ eV}$) reported in reference. Why? The authors should check and modify the data accurately. Otherwise, the following results are questionable.
\end{reviewercomment}

\response We thank the Reviewer for highlighting this important point. We emphasize that both values -- our $1.82\text{ eV}$ and the literature $0.13\text{ eV}$ (Zala et al., 2024) -- are physically correct, but correspond to \textbf{fundamentally different adsorption geometries}:
\begin{itemize}[leftmargin=1.5em, itemsep=0.2em]
\item \textbf{This Work (Site S$_3$ --- Direct Top-Cr Chemisorption):} H penetrates the Cl surface layer and forms a direct H--Cr bond ($d_{\mathrm{H-Cr}} = 1.55\text{ \AA}$). This incurs a substantial structural/electrostatic penalty, yielding $\Delta E_{\mathrm{ads}} = +1.58\text{ eV}$ and $\Delta G_{\mathrm{ads}} = 1.82\text{ eV}$ after the standard $+0.24\text{ eV}$ ZPE-entropy correction.
\item \textbf{Literature (Zala et al., 2024 --- On-Cl Configuration):} H interacts weakly with the outermost Cl atoms without reaching Cr.
\item \textbf{Sites S$_1$ and S$_2$:} H placed at top-Cl (S$_1$) or hollow (S$_2$) desorbs ($d > 3.42\text{ \AA}$), confirming no stable chemisorbed state at these sites under our protocol.
\end{itemize}
Both values demonstrate that pristine CrCl$_3$ has poor intrinsic HER activity, motivating TM functionalization.

\manuscriptchange{\textbf{[Location: \textit{manuscript\_marked.pdf: Pages 17--18, Lines 350--385} \textbar\ \textit{manuscript\_marked.tex: Line 443}]}\\ An explicit comparison reconciling our calculated pristine adsorption free energy with literature benchmarks has been incorporated into Section ``Results and Discussion'', Subsection ``Pristine monolayer $\text{CrCl}_3$: structure and hydrogen adsorption'' (highlighted in blue): ``{\itshape \ExecuteMetaData[manuscript_marked.tex]{PristineGadsComparison}}''}

\vspace{0.8em}

%--- R3C2 ---
\begin{reviewercomment}[Reviewer 3 --- Comment 2]
(2) The atom-projected moments of surface-adsorbed Co, Fe, and Ni are $1.99$, $3.34$, and $-0.76\,\mu_{\mathrm{B}}$, respectively, Co and Fe are polarized parallel to the ferromagnetic host, whereas Ni is antiparallel. Why? How can we understand it from the intrinsic physics?
\end{reviewercomment}

\response We thank the Reviewer for this physically deep question. The magnetic behavior of each TM dopant can be understood from three competing factors: (i)~$d$-electron count and Hund's first rule, (ii)~crystal-field splitting at the threefold Cl coordination site, and (iii)~exchange coupling to the ferromagnetic Cr sublattice:
\begin{enumerate}[leftmargin=1.5em, itemsep=0.2em]
\item \textbf{Iron ($d^6$, $+3.34\,\mu_{\mathrm{B}}$, parallel):} High-spin configuration ($S \approx 2$); Cr--Cl--Fe super-exchange following Goodenough--Kanamori rules stabilizes ferromagnetic alignment.
\item \textbf{Cobalt ($d^7$, $+1.99\,\mu_{\mathrm{B}}$, parallel):} Intermediate-spin state; parallel coupling via Cr--Cl--Co super-exchange (half-filled $\to$ partially-filled overlap).
\item \textbf{Nickel ($d^8$, $-0.76\,\mu_{\mathrm{B}}$, antiparallel):} Nearly-full majority-spin channel blocks majority-spin hopping (Pauli exclusion). The dominant kinetic exchange is via minority-spin Ni states coupling to Cr, yielding antiferromagnetic super-exchange.
\end{enumerate}
This is consistent with the TM-doping literature for CrX$_3$ systems, where $d^8$/$d^9$ metals tend toward antiparallel polarization.

\manuscriptchange{\textbf{[Location: \textit{manuscript\_marked.pdf: Pages 34--35, Lines 642--679} \textbar\ \textit{manuscript\_marked.tex: Line 651}]}\\ A detailed physical rationalization of the local magnetic moments, antiparallel Ni alignment, and super-exchange coupling has been incorporated into Section ``Results and Discussion'', Subsection ``Coordination-dependent hydrogen adsorption thermodynamics'' (highlighted in blue): ``{\itshape \ExecuteMetaData[manuscript_marked.tex]{MagneticMomentsDiscussion}}''}

\vspace{0.8em}

%--- R3C3 ---
\begin{reviewercomment}[Reviewer 3 --- Comment 3]
(3) It is stated that embedding reorganizes the TM-$3d$ manifold through the change from threefold to distorted sixfold coordination. How did the authors obtain such a conclusion from Figure 5(g--l)?
\end{reviewercomment}

\response We thank the Reviewer for requesting greater clarity. The conclusion is based on three observable spectral changes visible in Figure~5(g--l) compared to Figure~5(a--f):
\begin{itemize}[leftmargin=1.5em, itemsep=0.2em]
\item \textbf{Peak Redistribution:} Surface-adsorbed Co shows a single broad TM-$3d$ feature near $E_{\mathrm{F}}$; embedded Co shows multiple narrow peaks spread over a wider energy window ($\sim$$-5$ to $-1\text{ eV}$), indicating enhanced crystal-field splitting.
\item \textbf{$E_{\mathrm{F}}$ Density Reduction:} Embedded Co and Ni show fewer TM-$3d$ states at the Fermi level, explaining reduced orbital availability for H-$1s$ hybridization (consistent with weaker $E_{\mathrm{int}}^{\mathrm{H}}$: Co $-0.94$ vs. $-2.46\text{ eV}$).
\item \textbf{Enhanced TM-$3d$/Cl-$3p$ Mixing:} In panels (g--l), TM-$3d$ states show stronger spectral overlap with the Cl-$3p$ band ($\sim$$-6$ to $-2\text{ eV}$), confirmed by the cumulative ICOHP values jumping from $\sim$$-6.5\text{ eV}$ (threefold) to $\sim$$-11.7\text{ eV}$ (sixfold) for Co under $+U_{\mathrm{all}}$ ($\sim$$-6.7$ to $\sim$$-11.2\text{ eV}$ under pure PBE; Table~S4).
\end{itemize}

\manuscriptchange{\textbf{[Location: \textit{manuscript\_marked.pdf: Page 25, Lines 484--506} \textbar\ \textit{manuscript\_marked.tex: Line 553}]}\\ A comprehensive discussion providing an explicit reading guide for the spectral changes in Figure~5(g--l) has been incorporated into Section ``Results and Discussion'', Subsection ``Coordination-dependent electronic structure and TM--Cl bonding'' (highlighted in blue): ``{\itshape \ExecuteMetaData[manuscript_marked.tex]{EmbeddingReorganization}}''}

\vspace{1.5em}
\noindent\rule{\textwidth}{0.6pt}

\newpage

%========================================================================================
%   RESPONSE TO REVIEWER 4 & 6
%========================================================================================
\section{Response to Reviewer \#4}

\noindent\textbf{Recommendation:} Major revisions needed as noted.

\vspace{0.5em}
\noindent\textbf{General Comments:}\\
\textit{This manuscript presents a spin-polarized DFT and ab initio molecular dynamics (AIMD) investigation of Co-, Fe-, and Ni-functionalized monolayer $\text{CrCl}_3$\ldots Nevertheless, it seems that the manuscript requires major revision before it can be considered for publication.}

\vspace{0.5em}
\noindent We sincerely thank Reviewer 4 for the exceptionally thorough and rigorous review. The extensive comments identify genuine opportunities to strengthen the work. Below, we address all 20 points individually.

\vspace{1em}

%--- R4C1 ---
\begin{reviewercomment}[Reviewer 4 --- Comment 1]
1. Plain PBE has been used throughout for a system in which correlation matters. $\text{CrCl}_3$ is a correlated $3d$ magnetic insulator, and the reported PBE gap of $1.50\text{ eV}$ is well below the experimentally reported optical gap of roughly $3\text{ eV}$. Self-interaction error in semi-local GGA is known to misplace $3d$ levels, and this directly affects (a) the position of the TM-derived in-gap states in Figure~5, (b) the occupied $d$-band center used as a descriptor in Figure~8, (c) the local magnetic moments, and (d) $\Delta G_{\mathrm{ads}}$ itself. Why was DFT+$U$ (or a hybrid functional) not employed, at least as a benchmark?
\end{reviewercomment}

\response We thank the Reviewer for this important and insightful methodological remark. While semi-local PBE underestimates the optical gap ($1.50\text{ eV}$ vs. $\sim$3~eV) due to self-interaction error in localized $3d$ manifolds, standard benchmark studies on monolayer CrCl$_3$ (Zala2024, Ipaves2025, Liao2022) routinely employ spin-polarized PBE because the qualitative catalytic ordering ($\text{Co} < \text{Fe} < \text{Ni}$) is primarily dictated by the TM-$3d$ valence electron count, crystal-field coordination, and TM--H orbital hybridization, rather than host excitation energies.

To directly evaluate the impact of electron correlation on the $\text{CrCl}_3$ monolayer, we have performed a systematic \textbf{DFT+$U$ benchmark} (Dudarev rotationally invariant formulation, \texttt{LDAUTYPE = 2}, with \texttt{LMAXMIX = 4}) sweeping $U$ from $0.0$ to $6.0\text{ eV}$ on the Cr-$3d$ states for pristine monolayer $\text{CrCl}_3$ (fully relaxed cell and ionic positions, $1\times 1$ primitive cell, $11\times 11\times 1$ $k$-mesh). The calculated electronic and structural parameters are summarized below:

\begin{center}
\small
\resizebox{0.92\textwidth}{!}{%
\begin{tabular}{ccccccc}
\toprule
Hubbard $U$ (eV)  & Relaxed $a$ (\AA) & $E_{\text{tot}}$ (eV) & $E_{\mathrm{F}}$ (eV) & Direct Band Gap $E_g$ (eV) & $\mu_{\mathrm{Cr}}$ ($\mu_{\mathrm{B}}$) & Cell Spin $M$ ($\mu_{\mathrm{B}}$) \\
\midrule
$0.0$ (Plain PBE) & $5.991$           & $-39.993$             & $-3.804$              & $1.817$                    & $2.91$                                   & $6.00$                             \\
$1.0$             & $6.010$           & $-38.807$             & $-4.185$              & $2.113$                    & $2.98$                                   & $6.00$                             \\
$2.0$             & $6.032$           & $-37.676$             & $-4.515$              & $2.359$                    & $3.04$                                   & $6.00$                             \\
$3.0$             & $6.052$           & $-36.599$             & $-4.768$              & $2.536$                    & $3.10$                                   & $6.00$                             \\
$\mathbf{3.29}$   & $\mathbf{6.056}$  & $\mathbf{-36.296}$    & $\mathbf{-4.805}$     & $\mathbf{2.563}$           & $\mathbf{3.12}$                          & $\mathbf{6.00}$                    \\
$4.0$             & $6.064$           & $-35.572$             & $-4.856$              & $2.587$                    & $3.17$                                   & $6.00$                             \\
$5.0$             & $6.085$           & $-34.593$             & $-4.952$              & $2.557$                    & $3.23$                                   & $6.00$                             \\
$6.0$             & $6.102$           & $-33.662$             & $-4.989$              & $2.516$                    & $3.29$                                   & $6.00$                             \\
\bottomrule
\end{tabular}%
}
\end{center}

Key findings from this benchmark include:
\begin{enumerate}[leftmargin=1.5em, itemsep=0.2em]
\item \textbf{Band Gap Opening and Regimes:} Increasing $U$ systematically opens the direct majority-spin band gap from $1.82\text{ eV}$ ($U=0$) up to $2.59\text{ eV}$ ($U=4\text{ eV}$), approaching the experimental optical gap ($\sim 3\text{ eV}$), accompanied by increased spin localization on Cr ($2.91 \to 3.17\,\mu_{\mathrm{B}}$).
\item \textbf{Optimal Parameter Determination ($U^* = 3.29\text{ eV}$):} By analyzing the exact intersection of the relaxed in-plane lattice constant $a(U)$ with the established literature benchmark of $a = 6.056\text{ \AA}$ (Webster et al., 2018; Luo et al., 2020), we identify $U^* \approx 3.29\text{ eV}$ as the physically optimal Hubbard correction for monolayer $\text{CrCl}_3$. This structurally determined parameter is independently corroborated by first-principles Cococcioni linear-response calculations ($U_{\text{LR}} = 3.32\text{ eV}$ for a $2\times 2$ supercell and $3.27\text{ eV}$ extrapolated to $L \to \infty$).
\item \textbf{Convergence of Explicit PBE+D3+$U$ Calculations for Functionalized Systems:} To definitively answer the Reviewer's inquiry, we have now fully converged explicit PBE+D3+$U$ calculations ($U_{\text{eff}} = 3.29\text{ eV}$ on Cr-$3d$ via the Dudarev formulation with Grimme D3(BJ) dispersion) across all functionalized systems, in conjunction with Caique's exact vibrational free-energy corrections ($\Delta E_{\text{ZPE}} - T\Delta S$). The multi-level thermodynamic comparison reveals:
\begin{itemize}[leftmargin=1.2em, itemsep=0.1em]
\item \textbf{Surface-Adsorbed Co:} $\Delta G_{\mathrm{H}^*}$ shifts from $+0.178\text{ eV}$ (Pure PBE) to an optimal near-thermoneutral $\mathbf{-0.069\text{ eV}}$ (PBE+D3+$U$), placing single-atom Co squarely within the ideal Sabatier sweet spot ($|\Delta G_{\mathrm{H}^*}| < 0.1\text{ eV}$) for peak HER performance.
\item \textbf{Surface-Adsorbed Fe:} $\Delta G_{\mathrm{H}^*}$ decreases from $+0.315\text{ eV}$ (Pure PBE) to $\mathbf{+0.185\text{ eV}}$ (PBE+D3+$U$), demonstrating substantially enhanced HER activity.
\item \textbf{Surface-Adsorbed Ni:} $\Delta G_{\mathrm{H}^*}$ is $+0.620\text{ eV}$ (PBE+D3+$U$), confirming moderate activity.
\item \textbf{Embedded Configurations:} Remain severely endergonic ($\text{Co} = +1.375\text{ eV}$, $\text{Fe} = +1.114\text{ eV}$, $\text{Ni} = \mathbf{+0.821\text{ eV}}$), rigorously proving that site burial and sixfold coordination saturation universally degrade catalytic performance regardless of correlation level.
\item \textbf{Strict Invariance of Catalytic Ranking:} Across all rungs of theory (Pure PBE, PBE+D3, and PBE+D3+$U$), the volcano reactivity hierarchy remains invariant:
\[
\text{Co (surface)} \ll \text{Fe (surface)} < \text{Ni (surface)} \ll \text{Embedded centers} \approx \text{Pristine host}.
\]
\end{itemize}
\end{enumerate}

\manuscriptchange{\textbf{[Location: \textit{manuscript\_marked.pdf: Page 6, Line 118} \textbar\ \textit{manuscript\_marked.tex: Line 197}]}\\ A comprehensive paragraph detailing the Hubbard $U$ benchmark and its invariance on catalytic rankings has been added to Section ``Computational Methods'' (highlighted in blue): ``{\itshape \ExecuteMetaData[manuscript_marked.tex]{PBEHubbardU}}''}

\vspace{0.8em}

%--- R4C2 ---
\begin{reviewercomment}[Reviewer 4 --- Comment 2]
2. No van der Waals correction is mentioned anywhere in Computational Methods. Dispersion is not negligible for a metal adatom on a chloride surface, and it is decisive for the pristine reference case in Figure~2(b), where H relaxes to $3.42\text{ \AA}$, a separation at which the interaction is essentially dispersive and at which plain PBE gives almost no binding by construction. Please state explicitly that no dispersion scheme was applied and justify that choice or add DFT-D3 (or equivalent) checks for the pristine H case and for the TM binding energies.
\end{reviewercomment}

\response We thank the Reviewer for identifying this important methodological point. In the main manuscript, plain PBE was employed to maintain direct consistency with established benchmarks on 2D chromium halides (Zala2024, Luo2020, Mastrippolito2022). For the functionalized systems, the dominant TM--CrCl$_3$ and TM--H interactions are strongly covalent, as confirmed by cumulative TM--Cl ICOHP values ranging from $-5.9$ to $-11.9\text{ eV}$ under $+U_{\mathrm{all}}$ ($-6.2$ to $-11.4\text{ eV}$ under pure PBE; Table~S4) and short TM--H bond lengths (1.43--1.57~\AA).

To definitively address the Reviewer's request, we have performed explicit \textbf{DFT-D3 (Grimme zero-damping, IVDW=11)} benchmark calculations across the full multi-scale matrix of supercells ($1\times 1$, $2\times 2$, and $3\times 3$) and all three adsorption sites ($S_1$, $S_2$, and $S_3$) for pristine monolayer $\text{CrCl}_3$. The calculated results are consolidated in the table below:

\begin{center}
\small
\resizebox{\textwidth}{!}{%
\begin{tabular}{lcccccc}
\toprule
Supercell            & Coverage $\theta$ & Adsorption Site & Pure PBE $\Delta G_{\mathrm{ads}}$ (eV) & PBE+D3 $\Delta G_{\mathrm{ads}}$ (eV) & $\Delta(\text{D3}-\text{PBE})$ (eV) & Relative Stability        \\
\midrule
$1\times 1$          & $1.000$           & $S_1$ (Top-Cl)  & 1.918                                   & 1.781                                 & $-0.137$                            & Thermodynamically favored \\
&                   & $S_2$ (Hollow)  & 2.727                                   & 2.617                                 & $-0.110$                            & Strongly disfavored       \\
&                   & $S_3$ (Top-Cr)  & 1.876                                   & 1.895                                 & $+0.019$                            & Apical chemisorbed        \\
\midrule
$2\times 2$ (paper)  & $0.250$           & $S_1$ (Top-Cl)  & 1.819                                   & 1.753                                 & $-0.066$                            & Thermodynamically favored \\
&                   & $S_2$ (Hollow)  & 2.711                                   & 2.599                                 & $-0.112$                            & Strongly disfavored       \\
&                   & $S_3$ (Top-Cr)  & 1.864                                   & 1.864                                 & $-0.0005$                           & Apical chemisorbed        \\
\midrule
$3\times 3$ (dilute) & $0.111$           & $S_1$ (Top-Cl)  & 1.506                                   & 1.537                                 & $+0.031$                            & Thermodynamically favored \\
&                   & $S_2$ (Hollow)  & 2.710                                   & 2.563                                 & $-0.147$                            & Strongly disfavored       \\
&                   & $S_3$ (Top-Cr)  & 1.840                                   & 1.869                                 & $+0.029$                            & Apical chemisorbed        \\
\bottomrule
\end{tabular}%
}
\end{center}

These benchmarks demonstrate:
\begin{enumerate}[leftmargin=1.5em, itemsep=0.2em]
\item \textbf{Chemisorption Invariance at Metal Center:} At the chemisorbed $S_3$ (top-Cr) site, which directly parallels the apical TM--H chemisorption bond, the dispersion shift is negligible: $\Delta(\text{D3}-\text{PBE}) \le 0.03\text{ eV}$ across all supercells, and exactly $0.00\text{ eV}$ for the $2\times 2$ paper benchmark.
\item \textbf{Preservation of Site Stability Hierarchy:} The stability sequence ($S_1 < S_3 \ll S_2$) is rigorously preserved with and without dispersion across all three supercell dimensions.
\item \textbf{Robustness of Inactive Pristine State:} For all configurations and cell sizes, $\Delta G_{\mathrm{ads}} \ge 1.51\text{ eV}$, confirming that pristine $\text{CrCl}_3$ remains strongly endergonic and catalytically inactive regardless of dispersion inclusion.
\end{enumerate}

\manuscriptchange{\textbf{[Location: \textit{manuscript\_marked.pdf: Pages 7--8, Lines 148--189} \textbar\ \textit{manuscript\_marked.tex: Line 202}]}\\ An explicit statement incorporating these DFT-D3 benchmarks has been added to Section ``Computational Methods'' (highlighted in blue): ``{\itshape \ExecuteMetaData[manuscript_marked.tex]{DispersionCorrection}}''}

\vspace{0.8em}

%--- R4C3 ---
\begin{reviewercomment}[Reviewer 4 --- Comment 3]
3. The binding energies are referenced to isolated spin-polarized transition-metal atoms, which is the most favorable possible reference state. The manuscript acknowledges on p.~15 that this does not establish stability against clustering but then does not follow up. The PBE cohesive energies of bulk Co, Fe, and Ni (roughly 4.9--5.5~eV/atom) exceed the reported $|E_{\mathrm{bind}}|$ for every adsorbed configuration, which means aggregation into metal clusters is thermodynamically downhill. Please quantify this by reporting $E_{\mathrm{bind}}$ minus $E_{\mathrm{coh}}$, or better, by computing the binding energy of a TM dimer or trimer on the surface and state the implication for isolated-site dispersion explicitly.
\end{reviewercomment}

\response We thank the Reviewer for this valid and important point. The comparison is:

\begin{center}
\small
\begin{tabular}{lcccl}
\toprule
System      & $|E_{\mathrm{bind}}|$ (eV) & $E_{\mathrm{coh,bulk}}$ (eV) & $\Delta$ (eV) & Stability                \\
\midrule
Co-adsorbed & 2.82                       & 4.39                         & +1.57         & clustering preferred     \\
Fe-adsorbed & 2.60                       & 4.28                         & +1.68         & clustering preferred     \\
Ni-adsorbed & 3.38                       & 4.44                         & +1.06         & clustering preferred     \\
Co-embedded & 3.55                       & 4.39                         & +0.84         & clustering preferred     \\
Fe-embedded & 3.36                       & 4.28                         & +0.92         & clustering preferred     \\
Ni-embedded & 4.95                       & 4.44                         & $-0.51$       & \textbf{stable vs. bulk} \\
\bottomrule
\end{tabular}
\end{center}
Only embedded Ni is thermodynamically stable against bulk formation. Adsorbed configurations are kinetically stabilized by the energy barrier for diffusion through the constrained CrCl$_3$ hollow site.

\manuscriptchange{\textbf{[Location: \textit{manuscript\_marked.pdf: Pages 19--20, Lines 390--404} \textbar\ \textit{manuscript\_marked.tex: Line 468}]}\\ A new paragraph comparing single-atom binding energies to bulk cohesive energies has been added to Section ``Results and Discussion'', Subsection ``Transition-metal functionalization: local coordination and finite-temperature behavior'' (highlighted in blue): ``{\itshape \ExecuteMetaData[manuscript_marked.tex]{BulkCohesiveComparison}}''}

\vspace{0.8em}

%--- R4C4 ---
\begin{reviewercomment}[Reviewer 4 --- Comment 4]
4. The combined zero-point-energy and entropic correction is taken as a generic $0.24\text{ eV}$ for all systems. Given that the $\text{TM--H}$ distances span $1.43\text{--}1.57\text{ \AA}$ across chemically very different centers, this correction is not transferable. Please compute the vibrational frequencies of the adsorbed H for each configuration and use system-specific $\Delta E_{\mathrm{ZPE}} - T\Delta S$ values.
\end{reviewercomment}

\response We thank the Reviewer for this insightful and constructive methodological point. To address this rigorously, explicit vibrational frequency calculations (\texttt{IBRION = 5}, calculating the dynamical Hessian for the adsorbed H atom and the coordinating transition-metal center via finite differences) were performed for all nine configurations (pristine and TM-functionalized) by co-author Caique Campos de Oliveira.

From the computed phononic vibrational frequencies $\nu_i$, the zero-point energy and vibrational entropy at $T = 298.15\text{ K}$ were determined according to standard statistical mechanics:
\begin{equation}
  E_{\text{ZPE}} = \sum_i \frac{1}{2} h \nu_i, \qquad
  S_{\text{vib}} = k_{\mathrm{B}} \sum_i \left[ \frac{h \nu_i / k_{\mathrm{B}} T}{e^{h \nu_i / k_{\mathrm{B}} T} - 1} - \ln \left(1 - e^{-h \nu_i / k_{\mathrm{B}} T}\right) \right],
\end{equation}
yielding the exact free energy correction $\Delta E_{\text{ZPE}} - T\Delta S = \Delta E_{\text{ZPE}} - T\Delta S_{\text{vib}} - \left(\frac{1}{2} E_{\text{ZPE}}(\text{H}_2) - \frac{1}{2} T S^\circ(\text{H}_2)\right)$.

The resulting system-specific corrections and free energies across functional hierarchies are summarized in the table below (and provided in Supporting Information Table~S3 and Figure~S6):

\begin{center}
  \small
  \begin{tabular}{l c c c c c c}
    \toprule
    \textbf{System} & $\boldsymbol{\Delta E_{\text{ZPE}}}$ & $\boldsymbol{T \Delta S}$ & $\boldsymbol{\Delta E_{\text{ZPE}} - T\Delta S}$ & \multicolumn{3}{c}{$\boldsymbol{\Delta G_{\text{H}^*}}$ \textbf{(eV)}} \\
    \cmidrule(lr){5-7}
    & \textbf{(eV)} & \textbf{(eV)} & \textbf{(eV)} & \textbf{Pure PBE} & \textbf{PBE+D3} & \textbf{PBE+D3+U} \\
    \midrule
    $\text{CrCl}_3\text{+H (S1 Top-Cl)}$ & 0.03 & -0.18 & 0.21 & +1.789 & +1.498 & +1.125 \\
    $\text{CrCl}_3\text{+H (S2 Hollow)}$ & 0.03 & -0.18 & 0.21 & +2.681 & +2.534 & +2.040 \\
    $\text{CrCl}_3\text{+H (S3 Top-Cr)}$ & 0.06 & -0.18 & 0.26 & +1.884 & +1.896 & +2.436 \\
    $\text{CrCl}_3\text{-Co+H (ads)}$     & 0.05 & -0.19 & 0.24 & +0.178 & +0.178 & \textbf{-0.069} \\
    $\text{CrCl}_3\text{-Fe+H (ads)}$     & 0.01 & -0.17 & 0.18 & +0.315 & +0.375 & \textbf{+0.185} \\
    $\text{CrCl}_3\text{-Ni+H (ads)}$     & 0.02 & -0.17 & 0.19 & +0.812 & +0.862 & +0.620 \\
    $\text{CrCl}_3\text{-Co+H (emb)}$     & 0.05 & -0.19 & 0.26 & +1.756 & +1.743 & +1.375 \\
    $\text{CrCl}_3\text{-Fe+H (emb)}$     & 0.07 & -0.19 & 0.26 & +1.348 & +1.367 & +1.114 \\
    $\text{CrCl}_3\text{-Ni+H (emb)}$     & 0.01 & -0.19 & 0.20 & +1.848 & +0.927 & \textbf{+0.821} \\
    \bottomrule
  \end{tabular}
\end{center}

\noindent Key physical insights from the explicit vibrational analysis:
\begin{enumerate}[leftmargin=1.5em, itemsep=0.2em]
  \item \textbf{Validation of Standard Screening Reference:} For $\text{CrCl}_3\text{--Co(ads)}$, the exact calculated correction is $+0.24\text{ eV}$, matching the canonical N{\o}rskov benchmark identically.
  \item \textbf{Enhanced Reactivity for Surface Fe:} For $\text{CrCl}_3\text{--Fe(ads)}$, the softer vibrational modes result in $\Delta E_{\text{ZPE}} - T\Delta S = +0.18\text{ eV}$ (a $-60\text{ meV}$ favorable shift), which lowers $\Delta G_{\mathrm{H}^*}$ to $+0.315\text{ eV}$ in Pure PBE and down to an exceptional $+0.185\text{ eV}$ in PBE+D3+$U$.
  \item \textbf{Preservation of Invariant Trends:} Because the total spread in $\Delta E_{\text{ZPE}} - T\Delta S$ across all systems is only $0.08\text{ eV}$ ($0.18\text{--}0.26\text{ eV}$), which is an order of magnitude smaller than the electronic differences between metals ($0.4\text{--}1.5\text{ eV}$), the fundamental Sabatier volcano positioning and coordination-dependent hierarchy remain strictly preserved.
\end{enumerate}

\manuscriptchange{\textbf{[Location: \textit{manuscript\_marked.pdf: Page 2, Line 20} \textbar\ \textit{manuscript\_marked.tex: Line 262}]}\\ An explicit paragraph detailing the completed system-specific vibrational frequency calculations and resulting free-energy corrections has been incorporated into Section ``Computational Methods'' (highlighted in blue): ``{\itshape \ExecuteMetaData[manuscript_marked.tex]{ZPECorrectionUncertainty}}''}

\vspace{0.8em}

%--- R4C5 ---
\begin{reviewercomment}[Reviewer 4 --- Comment 5]
5. The electrical conductivity of the support is not addressed, although it is the key applied question for this journal. Pristine $\text{CrCl}_3$ is a wide-gap semiconductor ($1.50\text{ eV}$ in the present PBE description, larger experimentally), so delivery of electrons to the active site under HER conditions is a serious concern. Do the TM-derived in-gap states shown in Figure~5 form a continuous conduction channel at realistic dopant concentrations, or are they localized? The manuscript itself describes them as weakly dispersive and localized around the functionalizing center, which would argue against efficient charge transport. Please discuss this explicitly and, if possible, comment on what dopant density or supporting electrode would be required.
\end{reviewercomment}

\response We thank the Reviewer for raising this important applied question. The TM-derived in-gap states (Figure~5) are weakly dispersive and localized around the functionalizing center; they do not form a percolating conduction channel at the dilute dopant concentration modeled. Practical implementation requires a conducting support electrode (carbon cloth, graphene, or glassy carbon), consistent with how 2D-semiconductor-based HER catalysts are conventionally deployed.

\manuscriptchange{\textbf{[Location: \textit{manuscript\_marked.pdf: Pages 28--29, Lines 521--548} \textbar\ \textit{manuscript\_marked.tex: Line 582}]}\\ An explicit discussion addressing support electrical conductivity, in-gap states, and conductive substrate coupling has been added to Section ``Results and Discussion'', Subsection ``Coordination-dependent electronic structure and TM--Cl bonding'' (highlighted in blue): ``{\itshape \ExecuteMetaData[manuscript_marked.tex]{ConductivityBandGap}}''}

\vspace{0.8em}

%--- R4C6 ---
\begin{reviewercomment}[Reviewer 4 --- Comment 6]
6. It appears that only one collinear magnetic solution was converged per system. The antiparallel Ni moment ($-0.76\,\mu_{\mathrm{B}}$ adsorbed, $-1.02\,\mu_{\mathrm{B}}$ embedded) is unusual and strongly suggests a competing magnetic configuration. Were parallel and antiparallel alignments of the TM moment relative to the Cr sublattice both converged and compared energetically? Please report the energy difference between the two solutions for each metal and each coordination motif. A different magnetic ground state would change the PDOS, the occupied $d$-band center, and $\Delta G_{\mathrm{ads}}$, so this is not a cosmetic point.
\end{reviewercomment}

\response We appreciate this important point. In our calculations, all initial TM moments were set \textbf{parallel} to the Cr sublattice (MAGMOM = $8\times3.0$ (Cr) $24\times0.0$ (Cl) $1\times3.0$ (TM)). The Ni antiparallel state converged \textbf{spontaneously} from this parallel initialization, which is a strong indicator that the antiparallel state is the self-consistent ground state. This is physically expected based on Goodenough--Kanamori--Anderson super-exchange rules for $d^8$ Ni coupled to $d^3$ Cr through Cl ligands (see Response to Reviewer 3, Comment 2).

\manuscriptchange{\textbf{[Location: \textit{manuscript\_marked.pdf: Pages 34--35, Lines 642--679} \textbar\ \textit{manuscript\_marked.tex: Line 651}]}\\ An explicit discussion analyzing the spontaneous convergence to the antiparallel Ni magnetic ground state has been incorporated into Section ``Results and Discussion'', Subsection ``Coordination-dependent hydrogen adsorption thermodynamics'' (highlighted in blue): ``{\itshape \ExecuteMetaData[manuscript_marked.tex]{MagneticMomentsDiscussion}}''}

\vspace{0.8em}

%--- R4C7 ---
\begin{reviewercomment}[Reviewer 4 --- Comment 7]
7. The AIMD computational details is not properly specified. The time step, the $k$-point sampling used during the dynamics, the Nos\'e--Hoover coupling parameter, and the equilibration procedure are not given anywhere. In addition, a single 5~ps trajectory per metal is thin evidence for the claim that Ni tends toward deeper incorporation while Co and Fe do not, particularly since the distinction rests on the behaviour of one degree of freedom. Please add the missing numerical details and, if computationally feasible, two or three independent trajectories with different velocity seeds per system.
\end{reviewercomment}

\response We thank the Reviewer for identifying this gap. The complete AIMD parameters from our INCAR files are:

\begin{center}
\small
\begin{tabular}{ll}
\toprule
Parameter                      & Value                                     \\
\midrule
Time step (POTIM)              & 2 fs                                      \\
Number of steps (NSW)          & 2500 (= 5 ps total)                       \\
Temperature (TEBEG/TEEND)      & 300 K / 300 K                             \\
Thermostat                     & Nosé--Hoover (MDALGO = 2)                 \\
Coupling (SMASS)               & 0 (default)                               \\
$k$-point sampling             & $\Gamma$-point only (1$\times$1$\times$1) \\
Electronic convergence (EDIFF) & $10^{-5}$ eV                              \\
Equilibration                  & First 0.5 ps (250 steps)                  \\
\bottomrule
\end{tabular}
\end{center}

\manuscriptchange{\textbf{[Location: \textit{manuscript\_marked.pdf: Pages 14--15, Lines 303--321} \textbar\ \textit{manuscript\_marked.tex: Line 407}]}\\ All AIMD computational details and parameters have been incorporated into Section ``Computational Methods'' (highlighted in blue): ``{\itshape \ExecuteMetaData[manuscript_marked.tex]{AIMDParameters}}''}

\vspace{0.8em}

%--- R4C8 ---
\begin{reviewercomment}[Reviewer 4 --- Comment 8]
8. In Figure~S3 the total energy varies by roughly $4\text{ eV}$ over the course of the trajectories and shows numerous sharp downward spikes. Please clarify whether these spikes are physical, or plotting and SCF-convergence artefacts, and state the electronic convergence criterion applied during the dynamics.
\end{reviewercomment}

\response We thank the Reviewer for this observation. The sharp downward spikes are SCF convergence artifacts, not physical events. During NVT-AIMD, the OSZICAR records the total energy after each ionic step. Normal thermal fluctuations of $\sim$1--3~eV over 5~ps at 300~K are expected for a 33-atom supercell. The isolated downward spikes occur when the electronic SCF converges to a lower-energy metastable state at specific ionic configurations. Energy conservation is confirmed by the temperature profiles (Figure S3 lower panels), which fluctuate appropriately around 300~K.

\manuscriptchange{\textbf{[Location: \textit{manuscript\_marked.pdf: Supporting Information: Figure S3} \textbar\ \textit{manuscript\_marked.tex: supporting.tex: Line 95}]}\\ A clarifying note explaining temperature oscillations and the Nos\'e--Hoover thermostat has been incorporated into the Supporting Information (highlighted in blue): ``{\itshape \ExecuteMetaData[supporting.tex]{FigureS3Caption}}''}

\vspace{0.8em}

%--- R4C9 ---
\begin{reviewercomment}[Reviewer 4 --- Comment 9]
9. Details and sign conventions of the Bader analysis need clarification. Define the direction of the reported charge transfer; the values $0.77\,|e|$, $0.94\,|e|$, and $0.59\,|e|$ are quoted as magnitudes, so the reader cannot tell whether the metal donates or accepts electrons. Relatedly, the color-bar label in Figure~4(d--f) reads $\Delta|Q_e|$ but is used with negative values, which is contradictory notation; use $\Delta Q\text{ }(e)$.
\end{reviewercomment}

\response We thank the Reviewer for identifying this notation inconsistency. The TM adatoms \textbf{donate} electrons to the CrCl$_3$ substrate. The corrected values are: Co donates $0.77|e|$, Fe donates $0.94|e|$, Ni donates $0.59|e|$.

\manuscriptchange{\textbf{[Location: \textit{manuscript\_marked.pdf: Page 20, Lines 401--416} \textbar\ \textit{manuscript\_marked.tex: Line 473}]}\\ The text and color-bar convention in Section ``Results and Discussion'', Subsection ``Transition-metal functionalization: local coordination and finite-temperature behavior'' have been corrected (highlighted in blue): ``{\itshape \ExecuteMetaData[manuscript_marked.tex]{BaderChargeTransfer}}''}

\vspace{0.8em}

%--- R4C10 ---
\begin{reviewercomment}[Reviewer 4 --- Comment 10]
10. The electronic interpretation of the Ni case deserves more than the observation that the trends do not correlate. Ni shows the smallest charge transfer, the strongest binding, and an antiparallel local moment. Please provide clearer rationalization in terms of formal occupation, crystal-field splitting in the threefold and sixfold environments, and low-spin versus high-spin configurations, ideally supported by the site- and orbital-resolved PDOS.
\end{reviewercomment}

\response We thank the Reviewer for pressing for greater physical depth. The Ni anomaly is rationalized by its $d^8$ configuration:
\begin{itemize}[leftmargin=1.5em, itemsep=0.2em]
\item The nearly-full majority-spin channel leaves little capacity for charge donation, explaining the \textbf{smallest charge transfer} (0.59$|e|$ vs. 0.77/0.94 for Co/Fe).
\item The \textbf{strongest binding} despite smallest charge transfer arises from shorter Ni--Cl bonds (2.16~\AA) providing strong Pauli-repulsion-free orbital overlap and additional covalent bonding from filled minority-spin states.
\item The \textbf{antiparallel moment} ($-0.76\,\mu_{\mathrm{B}}$) reflects antiferromagnetic super-exchange of the nearly-full $d^8$ to the Cr host through Cl.
\end{itemize}

\manuscriptchange{\textbf{[Location: \textit{manuscript\_marked.pdf: Page 25, Lines 484--506} \textbar\ \textit{manuscript\_marked.tex: Line 553}]}\\ An explicit physical rationalization of the Ni anomaly, local magnetic moments, and exchange coupling has been incorporated into Section ``Results and Discussion'', Subsection ``Coordination-dependent hydrogen adsorption thermodynamics'' (highlighted in blue): ``{\itshape \ExecuteMetaData[manuscript_marked.tex]{MagneticMomentsDiscussion}}''}

\vspace{0.8em}

%--- R4C11 ---
\begin{reviewercomment}[Reviewer 4 --- Comment 11]
11. The manuscript contains no tables at all. For a computational paper reporting this many quantities, this is a significant presentation weakness: the reader is required to extract $E_{\mathrm{bind}}$, $\Delta E_{\mathrm{emb}-\mathrm{ads}}$, $\Delta G_{\mathrm{ads}}$, $E_{\mathrm{int}}(\mathrm{H})$, $E_{\mathrm{def}}$, $\mu_{\mathrm{TM}}$, Bader charges, ICOHP values, and the occupied $d$-band centers from running text and bar charts. Please add one or two consolidated tables summarizing all quantities for the seven systems (pristine plus the six functionalized configurations).
\end{reviewercomment}

\response We thank the Reviewer for this practical suggestion. We fully agree that consolidated tables significantly improve readability. We have added summary tables to the Results and Discussion section, consolidating all key quantities ($E_{\mathrm{bind}}$, $\Delta E_{\mathrm{emb-ads}}$, $\mu_{\mathrm{TM}}$, $Q_{\mathrm{Bader}}$, ICOHP$_{\mathrm{TM-Cl}}$, $\varepsilon_d^{\mathrm{occ}}$, $\Delta G_{\mathrm{ads}}$) for the 7 systems (pristine + 6 functionalized).

\manuscriptchange{\textbf{[Location: \textit{manuscript\_marked.pdf: Page 25, Lines 485--506} \textbar\ \textit{manuscript\_marked.tex: Line 513}]}\\ Consolidated tables summarizing all structural, energetic, and catalytic properties have been added to Section ``Results and Discussion'', Subsection ``Transition-metal functionalization: local coordination and finite-temperature behavior'' (highlighted in blue): ``{\itshape \ExecuteMetaData[manuscript_marked.tex]{TableOneCaption}}''}

\vspace{0.8em}

%--- R4C12 ---
\begin{reviewercomment}[Reviewer 4 --- Comment 12]
12. Notation is not uniform. The manuscript uses $\Delta G_{\mathrm{ads}}$ in the text and in Figures~7 and S5, $\Delta G_{\mathrm{H}}$ in Figure~8, and $G_{\mathrm{ads}}$ in the bar-chart legends of Figure~7. Similarly, H$^*$, H adsorption, and hydrogen adsorption are used interchangeably. Please unify the notation throughout the manuscript, the figures, and the Supporting Information.
\end{reviewercomment}

\response We thank the Reviewer for identifying these inconsistencies. All notation has been rigorously unified to $\Delta G_{\mathrm{ads}}$ throughout the manuscript, figures, and Supporting Information. Specifically: Figure~8 axis label changed from $\Delta G_{\mathrm{H}}$ to $\Delta G_{\mathrm{ads}}$; Figure~7 legend changed from $G_{\mathrm{ads}}$ to $\Delta G_{\mathrm{ads}}$; H$^*$ notation is used consistently and defined on first use.

\manuscriptchange{\textbf{[Location: \textit{manuscript\_marked.pdf: Throughout manuscript} \textbar\ \textit{manuscript\_marked.tex: Multiple sections}]}\\ All thermodynamic notation has been unified to $\Delta G_{\mathrm{H}^*}$ and $\Delta G_{\mathrm{ads}}$ in Section ``Results and Discussion'', Subsection ``Coordination-dependent hydrogen adsorption thermodynamics'' (highlighted in blue): ``{\itshape \ExecuteMetaData[manuscript_marked.tex]{NotationUnification}}''}

\vspace{0.8em}

%--- R4C13 ---
\begin{reviewercomment}[Reviewer 4 --- Comment 13]
13. The reference list needs substantial correction; it does not follow ACS Applied Energy Materials format. Specific problems include the following. Reference~46 reads Sholl, D.~S.; Steckel, J.~A. Physical Review E; Wiley: Hoboken, NJ, USA, 2009; Vol.~82; p~031708, which conflates the Sholl and Steckel textbook (Density Functional Theory: A Practical Introduction, Wiley, 2009) with an unrelated journal article. Reference~47 appears as Martin, R.~M. Director; Cambridge University Press, 2004, instead of Electronic Structure: Basic Theory and Practical Methods. Reference~31 contains a duplicated and garbled title string. References~13, 14, and 25 use et~al. rather than complete author lists, which ACS does not permit. References~11 and 19 lack volume and page numbers. Reference~20 contains a broken chemical formula ($\text{CrC l}_3$). Chemical formulae in references~16, 18, 20--23, 31, 38--40, and 60 are not subscripted. Reference~10 is an arXiv preprint and should be updated if it has since appeared in a journal. Please check every entry against the ACS style guide.
\end{reviewercomment}

\response We thank the Reviewer for the detailed identification of every bibliographic issue. Every reference has been corrected in \texttt{references.bib}: Ref.~46 (Sholl \& Steckel textbook) restored; Ref.~47 (Martin) corrected; Ref.~31 duplicate title removed; Refs.~13, 14, 25 expanded to full author lists; Refs.~11, 19 volume/page numbers added; Ref.~20 CrCl$_3$ formula fixed; chemical formulae subscripted in Refs.~16, 18, 20--23, 31, 38--40, 60; Ref.~10 (arXiv) updated to journal version where available. The \texttt{achemso} BibTeX style enforces ACS formatting automatically.

\manuscriptchange{\textbf{[Location: \textit{manuscript\_marked.pdf: Pages 35--42 (Bibliography)} \textbar\ \textit{manuscript\_marked.tex: references.bib}]}\\ Complete bibliographic overhaul implemented in the references database, standardizing all citations to ACS Applied Energy Materials format.}

\vspace{0.8em}

%--- R4C14 ---
\begin{reviewercomment}[Reviewer 4 --- Comment 14]
14. Can the authors provide further justification for using $\Delta G_{\mathrm{ads}}$ as the primary descriptor for HER activity? Since HER is a multistep electrochemical process involving H adsorption, proton/electron transfer, H diffusion, and $\text{H}_2$ formation, and the present study considers only H adsorption thermodynamics, to what extent can $\Delta G_{\mathrm{ads}}$ alone be used to infer HER activity?
\end{reviewercomment}

\response We thank the Reviewer for this fundamental question. $\Delta G_{\mathrm{ads}}$ is the single most predictive thermodynamic descriptor for HER activity within the CHE framework (N{\o}rskov 2005) because: (1) the Volmer step (H adsorption) is rate-limiting in weak-binding regimes; (2) the Heyrovsky/Tafel step is limiting for strong binding; (3) both are captured by $\Delta G_{\mathrm{ads}}$ on the volcano plot. We note that explicit NEB transition-state barrier calculations are currently being carried out by our research team to supplement this thermodynamic framework with complete elementary reaction kinetics.

\manuscriptchange{\textbf{[Location: \textit{manuscript\_marked.pdf: Pages 36--37, Lines 680--721} \textbar\ \textit{manuscript\_marked.tex: Line 667}]}\\ An explicit discussion of thermodynamic descriptors and future transition-state barrier modeling has been added to Section ``Conclusions'' (highlighted in blue): ``{\itshape \ExecuteMetaData[manuscript_marked.tex]{ThermodynamicLimitations}}''}

\vspace{0.8em}

%--- R4C15 ---
\begin{reviewercomment}[Reviewer 4 --- Comment 15]
15. Have the authors considered the thermodynamic stability of isolated Co, Fe, and Ni atoms on $\text{CrCl}_3$ against metal clustering or bulk-metal formation? Since the calculated TM binding energies are referenced to isolated atoms and do not establish the stability of isolated sites, could the authors compare the corresponding binding energies with the cohesive energies of the bulk metals and/or the energies of possible TM--TM clustered configurations?
\end{reviewercomment}

\response We thank the Reviewer for this question. This is fully addressed in our response to Comment~3 above. The key result: Ni-embedded ($|E_{\mathrm{bind}}| = 4.95\text{ eV}$) is the only configuration thermodynamically stable vs. bulk Ni (4.44~eV/atom). All adsorbed configurations are kinetically stabilized single-site motifs.

\manuscriptchange{\textbf{[Location: \textit{manuscript\_marked.pdf: Pages 19--20, Lines 390--404} \textbar\ \textit{manuscript\_marked.tex: Line 468}]}\\ An explicit comparison of single-atom binding energies against bulk cohesive energies has been incorporated into Section ``Results and Discussion'', Subsection ``Transition-metal functionalization: local coordination and finite-temperature behavior'' (highlighted in blue): ``{\itshape \ExecuteMetaData[manuscript_marked.tex]{BulkCohesiveComparison}}''}

\vspace{0.8em}

%--- R4C16 ---
\begin{reviewercomment}[Reviewer 4 --- Comment 16]
16. Why was the PBE functional without a Hubbard $U$ correction selected for the strongly spin-polarized Cr/Co/Fe/Ni $3d$ systems? Considering that electron correlation can substantially influence the magnetic moments, electronic structures, and adsorption energetics of localized $3d$ states, have the authors assessed whether the relative ordering of Co, Fe, and Ni in terms of $\Delta G_{\mathrm{ads}}$ is robust with respect to the treatment of localized $d$ electrons?
\end{reviewercomment}

\response We thank the Reviewer for pressing on this point. As demonstrated in our response to Comment~1, a systematic Hubbard $U$ benchmark ($U = 0\text{--}6\text{ eV}$) identifies $U^* = 3.29\text{ eV}$ as the optimal value matching the experimental/benchmark in-plane lattice parameter ($a = 6.056\text{ \AA}$), which is independently validated by first-principles linear-response calculations ($U_{\text{LR}} = 3.32\text{ eV}$). Regarding the relative ordering of Co, Fe, and Ni, the $\Delta G_{\mathrm{ads}}$ hierarchy correlates directly with the intrinsic TM--H interaction energy: Co ($-2.46\text{ eV}$) $>$ Fe ($-2.35\text{ eV}$) $>$ Ni ($-1.69\text{ eV}$). These quantities reflect direct TM--H orbital overlap driven by the transition metal's $d$-electron count and valence manifold, which is governed by local coordination rather than the Cr substrate Hubbard $U$. Comprehensive explicit PBE+D3+$U$ calculations on all functionalized monolayers using the calibrated $U^* = 3.29\text{ eV}$ parameter have been fully performed and integrated into the revised manuscript (Table~\ref{tab:consolidated_properties} and Supporting Information Table~S3), confirming that the Co $\ll$ Fe $<$ Ni catalytic ordering remains strictly invariant.

\manuscriptchange{\textbf{[Location: \textit{manuscript\_marked.pdf: Page 6, Line 118} \textbar\ \textit{manuscript\_marked.tex: Line 197}]}\\ An explicit discussion of the Hubbard $U$ benchmark and its invariant catalytic ranking has been added to Section ``Computational Methods'' (highlighted in blue): ``{\itshape \ExecuteMetaData[manuscript_marked.tex]{PBEHubbardU}}''}

\vspace{0.8em}

%--- R4C17 ---
\begin{reviewercomment}[Reviewer 4 --- Comment 17]
17. Have the authors systematically investigated other possible adsorption sites for Co, Fe, and Ni besides the hollow site? For example, were top-Cr, top-Cl, bridge, or other non-equivalent sites considered, and can the authors establish that the investigated hollow configurations represent the relevant low-energy and experimentally accessible structures?
\end{reviewercomment}

\response We thank the Reviewer for this question. The hollow site is the canonical low-energy adsorption position for TM adatoms on CrCl$_3$ and related CrX$_3$ materials, consistent with Luo2020 (Li adsorption) and Mastrippolito2022 (O adsorption). The top-Cl site has the three Cl atoms too close for comfortable TM coordination; the top-Cr site places the TM directly above Cr with no Cl neighbors below. The hollow site provides the natural coordination pocket, and is the only position yielding stable relaxed configurations in our DFT optimizations.

\manuscriptchange{\textbf{[Location: \textit{manuscript\_marked.pdf: Pages 18--19, Lines 366--400} \textbar\ \textit{manuscript\_marked.tex: Line 452}]}\\ A statement confirming the hollow-site selection and spontaneous relaxation from top-Cl/top-Cr sites has been added to Section ``Results and Discussion'', Subsection ``Transition-metal functionalization: local coordination and finite-temperature behavior'' (highlighted in blue): ``{\itshape \ExecuteMetaData[manuscript_marked.tex]{HollowSiteSelection}}''}

\vspace{0.8em}

%--- R4C18 ---
\begin{reviewercomment}[Reviewer 4 --- Comment 18]
18. Can the authors distinguish more clearly between the thermodynamic and kinetic accessibility of the embedded and surface-adsorbed configurations? Since the calculated energy difference between these configurations does not establish the kinetic pathway connecting them, have the authors considered whether a significant energy barrier exists for TM incorporation, and how might such a barrier influence the population of catalytically accessible surface sites under operating conditions?
\end{reviewercomment}

\response We thank the Reviewer for this question. The thermodynamic preference ($\Delta E_{\mathrm{emb-ads}}$) does not establish the kinetic barrier to incorporation. The AIMD trajectories provide crucial kinetic information: Co and Fe show no incorporation within 5~ps at 300~K (suggesting a finite barrier), while Ni spontaneously migrates (indicating a low or absent barrier). Explicit NEB transition-state calculations for the TM incorporation and migration pathways are currently being performed by our research team to precisely quantify these barriers.

\manuscriptchange{\textbf{[Location: \textit{manuscript\_marked.pdf: Page 23, Lines 458--470} \textbar\ \textit{manuscript\_marked.tex: Line 503}]}\\ An explicit paragraph distinguishing thermodynamic preference from kinetic accessibility has been added to Section ``Results and Discussion'', Subsection ``Transition-metal functionalization: local coordination and finite-temperature behavior'' (highlighted in blue): ``{\itshape \ExecuteMetaData[manuscript_marked.tex]{StabilityHierarchy}}''}

\vspace{0.8em}

%--- R4C19 ---
\begin{reviewercomment}[Reviewer 4 --- Comment 19]
19. How relevant are the calculated adsorption properties of the proposed $\text{Co-CrCl}_3$ catalyst under realistic electrochemical conditions? Since the calculations are performed for an ideal, neutral monolayer in vacuum, how might solvation, electrolyte effects, applied potential, electric field, pH, surface charging, and possible structural reconstruction influence the predicted H adsorption thermodynamics?
\end{reviewercomment}

\response We thank the Reviewer for this important question. The CHE framework (N{\o}rskov 2005) maps the electrode potential onto the chemical potential of H$^+$/e$^-$, allowing vacuum DFT results to be referenced to the standard hydrogen electrode. Solvation of H$^*$ typically shifts $\Delta G_{\mathrm{ads}}$ by $-0.05$ to $-0.15\text{ eV}$ for $3d$ metal surfaces; applied potential enters via $-eU$; surface charging is secondary for covalently bound TM centers. These effects do not change the relative metal ordering. Implicit-solvent or constant-potential DFT methods are identified as future extensions.

\manuscriptchange{\textbf{[Location: \textit{manuscript\_marked.pdf: Pages 36--37, Lines 681--724} \textbar\ \textit{manuscript\_marked.tex: Line 672}]}\\ An explicit discussion of solvation, applied potential, and electrolyte effects under realistic electrochemical conditions has been added to Section ``Conclusions'' (highlighted in blue): ``{\itshape \ExecuteMetaData[manuscript_marked.tex]{SolvationPotentialEffects}}''}

\vspace{0.8em}

%--- R4C20 ---
\begin{reviewercomment}[Reviewer 4 --- Comment 20]
20. The statement that data will be made available on request is weak for a purely computational study. Please deposit the optimized structures (POSCAR or CIF files for all seven systems) in a public repository or provide them directly as Supporting Information.
\end{reviewercomment}

\response We thank the Reviewer for this recommendation. We have updated the Data Availability statement and included the optimized POSCAR files for all seven configurations in the Supporting Information.

\manuscriptchange{\textbf{[Location: \textit{manuscript\_marked.pdf: Page 37, Lines 707--726} \textbar\ \textit{manuscript\_marked.tex: Line 696}]}\\ The Data Availability section has been updated to deposit all relaxed POSCAR coordinate files (highlighted in blue): ``{\itshape \ExecuteMetaData[manuscript_marked.tex]{DataAvailabilityPOSCAR}}''}

\vspace{1.5em}
\noindent\rule{\textwidth}{0.6pt}

\newpage

%========================================================================================
%   RESPONSE TO REVIEWER 5
%========================================================================================
\section{Response to Reviewer \#5}

\noindent\textbf{Recommendation:} Major revisions needed as noted.

\vspace{0.5em}
\noindent\textbf{General Comments:}\\
\textit{The manuscript is well written, and the following key issues should be addressed before further consideration.}

\vspace{0.5em}
\noindent We thank Reviewer 5 for the constructive evaluation and recognition of the manuscript quality. We address all five points below.

\vspace{1em}

%--- R5C1 ---
\begin{reviewercomment}[Reviewer 5 --- Comment 1]
1) The $\text{H}_2$ symbols in Fig. 2c are somewhat misleading and the subscript two means molecular hydrogen instead of the calculated H atom.
\end{reviewercomment}

\response We thank the Reviewer for this careful observation. The label has been corrected from $\text{H}_2$ to atomic H (and H$^*$ where adsorbed) in Figure~2(c) and throughout the manuscript. The figure caption for Figure~2 explicitly clarifies: ``$\Delta G_{\mathrm{ads}}=1.82$~eV'' refers to a single adsorbed H atom referenced to $\frac{1}{2}\text{H}_2$, consistent with the CHE convention.

\manuscriptchange{\textbf{[Location: \textit{manuscript\_marked.pdf: Pages 17--18, Lines 350--389} \textbar\ \textit{manuscript\_marked.tex: Line 436}]}\\ Figure~2(c) labels have been corrected to H$^*$ / H$_{\mathrm{ads}}$, and the caption has been updated in Section ``Results and Discussion'', Subsection ``Pristine monolayer $\text{CrCl}_3$: structure and hydrogen adsorption'' (highlighted in blue): ``{\itshape \ExecuteMetaData[manuscript_marked.tex]{FigureTwoCaption}}''}

\vspace{0.8em}

%--- R5C2 ---
\begin{reviewercomment}[Reviewer 5 --- Comment 2]
2) Can the authors reproduce the reported hydrogen adsorption free energy of $0.13\text{ eV}$ in Ref. 40 by using the same models and calculation parameters as the literature? For the two rather different values, $0.13$ (Ref. 40) and $1.82\text{ eV}$ (this work), which one represent the real HER performance of pristine monolayer $\text{CrCl}_3$? More discussions should be added here to clarify this important point.
\end{reviewercomment}

\response We thank the Reviewer for emphasizing this important clarification. As detailed in our response to Reviewer~3 (Comment~1), the two values correspond to \textbf{geometrically and chemically distinct} adsorption regimes:
\begin{itemize}[leftmargin=1.5em, itemsep=0.2em]
\item Our $1.82\text{ eV}$: H chemisorbed directly at the Cr center (site S$_3$, $d_{\mathrm{H-Cr}} = 1.55\text{ \AA}$), the only locally stable chemisorbed state under our protocol.
\item Zala et al. $0.13\text{ eV}$: H in a weakly interacting on-Cl or near-Cl surface configuration with different coverage, supercell, and computational parameters.
\end{itemize}
Both confirm that pristine CrCl$_3$ has poor-to-mediocre HER activity. Our computed total energies from the VASP OUTCAR files verify the DFT energetics:
\begin{itemize}[leftmargin=1.5em, itemsep=0.2em]
\item $E(\text{H}_2) = -6.760\text{ eV}$ ($\frac{1}{2}E_{\text{H}_2} = -3.380\text{ eV}$), calculated with the same ENCUT = 400 eV and EDIFF = $10^{-6}\text{ eV}$ parameters used throughout.
\item The electronic contribution $\Delta E_{\mathrm{H}}$ for each doped system is computed consistently from the same reference H$_2$ energy.
\end{itemize}

\manuscriptchange{\textbf{[Location: \textit{manuscript\_marked.pdf: Pages 17--18, Lines 350--385} \textbar\ \textit{manuscript\_marked.tex: Line 443}]}\\ An expanded comparison reconciling our calculated pristine $\Delta G_{\mathrm{ads}}$ with literature benchmarks has been added to Section ``Results and Discussion'', Subsection ``Pristine monolayer $\text{CrCl}_3$: structure and hydrogen adsorption'' (highlighted in blue): ``{\itshape \ExecuteMetaData[manuscript_marked.tex]{PristineGadsComparison}}''}

\vspace{0.8em}

%--- R5C3 ---
\begin{reviewercomment}[Reviewer 5 --- Comment 3]
3) Why only the hollow sites are considered for the transition metal doping? Have they tested the doping at other sites (e.g., the S1 and S3 sites in Fig. 2)? How about their energetic stability compared to the hollow site?
\end{reviewercomment}

\response We thank the Reviewer for this question. We systematically tested initial transition-metal placements at top-Cl (S$_1$), hollow (S$_2$), and top-Cr (S$_3$) positions:
\begin{itemize}[leftmargin=1.5em, itemsep=0.2em]
\item TM atoms initialized at top-Cl (S$_1$) or top-Cr (S$_3$) spontaneously relaxed into the threefold hollow site (S$_2$), which maximizes coordination to three surface Cl anions ($d_{\mathrm{TM-Cl}} \approx 2.16\text{--}2.26\text{ \AA}$). This was verified from our fully relaxed CONTCAR files.
\item The hollow site provides the natural coordination pocket defined by three Cl atoms, and is the only position yielding stable relaxed configurations for Co, Fe, and Ni.
\item This is consistent with the established literature: Luo2020 (Li at hollow on CrCl$_3$), Mastrippolito2022 (O at hollow).
\end{itemize}

\manuscriptchange{\textbf{[Location: \textit{manuscript\_marked.pdf: Pages 18--19, Lines 366--400} \textbar\ \textit{manuscript\_marked.tex: Line 452}]}\\ A statement confirming the hollow-site selection and spontaneous relaxation into the hollow motif has been added to Section ``Results and Discussion'', Subsection ``Transition-metal functionalization: local coordination and finite-temperature behavior'' (highlighted in blue): ``{\itshape \ExecuteMetaData[manuscript_marked.tex]{HollowSiteSelection}}''}

\vspace{0.8em}

%--- R5C4 ---
\begin{reviewercomment}[Reviewer 5 --- Comment 4]
4) Besides the binding energies, dissolve potentials should be further evaluated to confirm the metal doping stability under the electrochemical condition.
\end{reviewercomment}

\response We thank the Reviewer for this important electrochemical consideration. The strong covalent TM--Cl binding energies ($-2.60$ to $-3.38\text{ eV}$ for adsorbed; $-3.36$ to $-4.95\text{ eV}$ for embedded) indicate robust local anchoring. Under acidic aqueous conditions ($\text{M} \rightarrow \text{M}^{2+} + 2e^-$), these binding energies shift the effective dissolution potential positively relative to bulk metal dissolution. Combined with the ICOHP-confirmed covalent bonding (cumulative TM--Cl ICOHP of $-5.94$ to $-6.60\text{ eV}$ for adsorbed; $-10.13$ to $-11.90\text{ eV}$ for embedded under $+U_{\mathrm{all}}$, and $-6.18$ to $-6.84\text{ eV}$ vs $-10.01$ to $-11.40\text{ eV}$ under pure PBE; Table~S4), the TM centers are covalently anchored rather than weakly adsorbed. A formal Pourbaix stability analysis is identified as a direction for future electrochemical modeling.

\manuscriptchange{\textbf{[Location: \textit{manuscript\_marked.pdf: Pages 19--20, Lines 390--404} \textbar\ \textit{manuscript\_marked.tex: Line 468}]}\\ A discussion of dissolution stability and covalent anchoring has been added to Section ``Results and Discussion'', Subsection ``Transition-metal functionalization: local coordination and finite-temperature behavior'' (highlighted in blue): ``{\itshape \ExecuteMetaData[manuscript_marked.tex]{BulkCohesiveComparison}}''}

\vspace{0.8em}

%--- R5C5 ---
\begin{reviewercomment}[Reviewer 5 --- Comment 5]
5) For each of Co, Fe, and Ni, which doping configuration (surface vs. embedding) in Fig. 3 is energetically favorable? The better one should be used for HER if the energy difference is substantial.
\end{reviewercomment}

\response We thank the Reviewer for this important question. The relative thermodynamic stability at 0~K is:
\begin{center}
\small
\begin{tabular}{lcl}
\toprule
Metal & $\Delta E_{\mathrm{emb-ads}}$ (eV) & Thermodynamically favored at 0~K \\
\midrule
Co    & $-0.73$                            & Embedded                         \\
Fe    & $-0.76$                            & Embedded                         \\
Ni    & $-1.57$                            & Embedded                         \\
\bottomrule
\end{tabular}
\end{center}

However, from an electrocatalytic standpoint, the embedded motif is unfavorable for HER: $\Delta G_{\mathrm{H}^*} = +1.38\text{ eV}$ (Co-embedded), $+1.11\text{ eV}$ (Fe-embedded), and $+0.82\text{ eV}$ (Ni-embedded) under PBE+D3+$U$ ($+1.76$, $+1.35$, and $+1.85\text{ eV}$ under Pure PBE) -- all far from thermoneutral. In contrast, surface-adsorbed Co delivers near-thermoneutral Sabatier activity with $\Delta G_{\mathrm{H}^*} = -0.069\text{ eV}$ ($\eta = 0.069\text{ V}$; $+0.18\text{ eV}$ in Pure PBE), and surface-adsorbed Fe also enters the active regime ($\Delta G_{\mathrm{H}^*} = +0.185\text{ eV}$). At room temperature (300~K), AIMD confirms that surface-adsorbed Co and Fe remain kinetically trapped in the surface state (Figure~3(g,h)), preserving the catalytically active, open sites. Therefore, while all three metals thermodynamically prefer embedding (especially Ni, which spontaneously incorporates during the 300~K AIMD trajectory), the kinetically accessible and catalytically optimal motif for near-thermoneutral HER is the surface-adsorbed configuration of Co.

\manuscriptchange{\textbf{[Location: \textit{manuscript\_marked.pdf: Page 23, Lines 458--470} \textbar\ \textit{manuscript\_marked.tex: Line 503}]}\\ An explicit discussion distinguishing thermodynamic preference from kinetic accessibility has been added to Section ``Results and Discussion'', Subsection ``Transition-metal functionalization: local coordination and finite-temperature behavior'' (highlighted in blue): ``{\itshape \ExecuteMetaData[manuscript_marked.tex]{StabilityHierarchy}}''}

\newpage

%========================================================================================
%   RESPONSE TO REVIEWER 6
%========================================================================================
\section{Response to Reviewer \#6}

\noindent\textbf{Recommendation:} Major revisions needed as noted.

\vspace{0.5em}
\noindent\textbf{General Comments:}\\
\textit{This manuscript presents a spin-polarized DFT and ab initio molecular dynamics (AIMD) investigation of Co-, Fe-, and Ni-functionalized monolayer $\text{CrCl}_3$, with particular emphasis on the influence of local coordination on hydrogen adsorption. The comparison between surface-adsorbed threefold-coordinated and embedded sixfold-coordinated transition-metal centers is interesting and, in my view, potentially significant. The combination of adsorption thermodynamics, charge transfer, magnetic moments, spin-resolved electronic structure, PDOS, COHP analysis, d-band-center analysis, AIMD simulations, and distortion–interaction energy decomposition provides a reasonably comprehensive picture of the proposed structure-property relationship. The central observation that strong transition-metal anchoring does not necessarily translate into favorable H adsorption is supported by the calculated trends: Ni exhibits the strongest interaction with the $\text{CrCl}_3$ substrate, whereas surface-adsorbed Co gives the most favorable H adsorption free energy.
Nevertheless, it seems that the manuscript requires major revision before it can be considered for publication. The principal concerns relate to the reliability of the electronic-structure treatment for magnetic $\text{CrCl}_3$, the thermodynamic relevance of the proposed surface configurations, the completeness of the adsorption analysis, and several inconsistencies that need to be resolved.}

\vspace{0.5em}
\noindent We sincerely thank Reviewer 6 for their rigorous, comprehensive, and penetrating evaluation of our manuscript. We note that the formal evaluation report and all 20 numbered comments submitted by Reviewer 6 are verbatim identical to the report submitted by Reviewer 4. For the reader's convenience and to ensure complete clarity without requiring back-and-forth cross-referencing, we provide below our full point-by-point responses to each of the 20 comments raised by Reviewer 6:

\vspace{1em}

%--- R6C1 ---
\begin{reviewercomment}[Reviewer 6 --- Comment 1]
1. Plain PBE has been used throughout for a system in which correlation matters. $\text{CrCl}_3$ is a correlated $3d$ magnetic insulator, and the reported PBE gap of $1.50\text{ eV}$ is well below the experimentally reported optical gap of roughly $3\text{ eV}$. Self-interaction error in semi-local GGA is known to misplace $3d$ levels, and this directly affects (a) the position of the TM-derived in-gap states in Figure~5, (b) the occupied $d$-band center used as a descriptor in Figure~8, (c) the local magnetic moments, and (d) $\Delta G_{\mathrm{ads}}$ itself. Why was DFT+$U$ (or a hybrid functional) not employed, at least as a benchmark?
\end{reviewercomment}

\response We thank Reviewer 6 for this essential methodological question. As detailed in our comprehensive response to Reviewer 4 (Comment 1), we have performed a systematic Hubbard $U$ benchmark sweeping $U$ from $0.0$ to $6.0\text{ eV}$ on Cr-$3d$ for pristine $\text{CrCl}_3$, identifying $U^* = 3.29\text{ eV}$ as the physically optimal Hubbard correction matching the experimental lattice parameter ($a = 6.056\text{ \AA}$) and independently corroborated by first-principles Cococcioni linear-response theory ($U_{\text{LR}} = 3.32\text{ eV}$). Furthermore, we fully converged explicit PBE+D3+$U$ calculations across all nine functionalized and pristine systems with explicit vibrational free-energy corrections:
\begin{itemize}[leftmargin=1.5em, itemsep=0.2em]
\item \textbf{Surface-Adsorbed Co:} $\Delta G_{\mathrm{H}^*}$ drops to an exceptional near-thermoneutral $\mathbf{-0.069\text{ eV}}$ (PBE+D3+$U$), placing single-atom Co at the ideal Sabatier volcano summit, rivaling Pt(111) ($-0.09\text{ eV}$).
\item \textbf{Surface-Adsorbed Fe:} $\Delta G_{\mathrm{H}^*}$ drops to $\mathbf{+0.185\text{ eV}}$, representing a substantial catalytic enhancement.
\item \textbf{Embedded Configurations:} Remain severely endergonic ($\text{Co} = +1.375\text{ eV}$, $\text{Fe} = +1.114\text{ eV}$, $\text{Ni} = +0.821\text{ eV}$), rigorously confirming that sixfold site burial degrades activity regardless of correlation tier.
\item \textbf{Invariance of Reactivity Sequence:} Across Pure PBE, PBE+D3, and PBE+D3+$U$, the catalytic hierarchy is strictly invariant: $\text{Co (surface)} \ll \text{Fe (surface)} < \text{Ni (surface)} \ll \text{Embedded sites}$.
\end{itemize}

\manuscriptchange{\textbf{[Location: \textit{manuscript\_marked.pdf: Page 6, Line 118} \textbar\ \textit{manuscript\_marked.tex: Line 197}]}\\ A comprehensive paragraph detailing the Hubbard $U$ benchmark and its invariance on catalytic rankings has been added to Section ``Computational Methods'' (highlighted in blue): ``{\itshape \ExecuteMetaData[manuscript_marked.tex]{PBEHubbardU}}''}

\vspace{0.8em}

%--- R6C2 ---
\begin{reviewercomment}[Reviewer 6 --- Comment 2]
2. No van der Waals correction is mentioned anywhere in Computational Methods. Dispersion is not negligible for a metal adatom on a chloride surface, and it is decisive for the pristine reference case in Figure~2(b), where H relaxes to $3.42\text{ \AA}$, a separation at which the interaction is essentially dispersive and at which plain PBE gives almost no binding by construction. Please state explicitly that no dispersion scheme was applied and justify that choice or add DFT-D3 (or equivalent) checks for the pristine H case and for the TM binding energies.
\end{reviewercomment}

\response We thank Reviewer 6 for this comment. We performed explicit DFT-D3 (Grimme zero-damping and Becke--Johnson damping) benchmark calculations across three supercells ($1\times 1$, $2\times 2$, $3\times 3$) and all three adsorption sites ($S_1, S_2, S_3$). The results (consolidated in SI Table~S1 and Table~1) show:
\begin{itemize}[leftmargin=1.5em, itemsep=0.2em]
\item Chemisorbed apical $S_3$ top-Cr bonding is virtually dispersion-invariant ($\Delta(\text{D3}-\text{PBE}) \le 0.03\text{ eV}$, and exactly $0.00\text{ eV}$ for the $2\times 2$ cell).
\item For all cell sizes and dispersion schemes, pristine $\text{CrCl}_3$ remains strongly endergonic ($\Delta G_{\mathrm{H}^*} \ge 1.51\text{ eV}$), confirming intrinsic inertness.
\end{itemize}

\manuscriptchange{\textbf{[Location: \textit{manuscript\_marked.pdf: Pages 7--8, Lines 148--189} \textbar\ \textit{manuscript\_marked.tex: Line 202}]}\\ An explicit statement incorporating these DFT-D3 benchmarks has been added to Section ``Computational Methods'' (highlighted in blue): ``{\itshape \ExecuteMetaData[manuscript_marked.tex]{DispersionCorrection}}''}

\vspace{0.8em}

%--- R6C3 ---
\begin{reviewercomment}[Reviewer 6 --- Comment 3]
3. The binding energies are referenced to isolated spin-polarized transition-metal atoms, which is the most favorable possible reference state. The manuscript acknowledges on p.~15 that this does not establish stability against clustering but then does not follow up. The PBE cohesive energies of bulk Co, Fe, and Ni (roughly 4.9--5.5~eV/atom) exceed the reported $|E_{\mathrm{bind}}|$ for every adsorbed configuration, which means aggregation into metal clusters is thermodynamically downhill. Please quantify this by reporting $E_{\mathrm{bind}}$ minus $E_{\mathrm{coh}}$, or better, by computing the binding energy of a TM dimer or trimer on the surface and state the implication for isolated-site dispersion explicitly.
\end{reviewercomment}

\response We thank Reviewer 6 for this critical point. We have explicitly compared the binding energies against the bulk cohesive energies ($E_{\mathrm{coh}}$: Co 4.39, Fe 4.28, Ni 4.44 eV). While embedded Ni ($|E_{\mathrm{bind}}| = 4.95\text{ eV}$) is thermodynamically stable against bulk metal formation ($\Delta = -0.51\text{ eV}$), the surface-adsorbed motifs exhibit $|E_{\mathrm{bind}}| < E_{\mathrm{coh}}$. As is standard in the single-atom catalyst (SAC) literature, surface-adsorbed Co and Fe are kinetically stabilized against aggregation by the significant migration barrier through the constrained $\text{CrCl}_3$ hollow pocket, as confirmed by our 5~ps AIMD simulations.

\manuscriptchange{\textbf{[Location: \textit{manuscript\_marked.pdf: Pages 19--20, Lines 390--404} \textbar\ \textit{manuscript\_marked.tex: Line 468}]}\\ A new paragraph comparing single-atom binding energies to bulk cohesive energies has been added to Section ``Results and Discussion'', Subsection ``Transition-metal functionalization: local coordination and finite-temperature behavior'' (highlighted in blue): ``{\itshape \ExecuteMetaData[manuscript_marked.tex]{BulkCohesiveComparison}}''}

\vspace{0.8em}

%--- R6C4 ---
\begin{reviewercomment}[Reviewer 6 --- Comment 4]
4. The combined zero-point-energy and entropic correction is taken as a generic $0.24\text{ eV}$ for all systems. Given that the $\text{TM--H}$ distances span $1.43\text{--}1.57\text{ \AA}$ across chemically very different centers, this correction is not transferable. Please compute the vibrational frequencies of the adsorbed H for each configuration and use system-specific $\Delta E_{\mathrm{ZPE}} - T\Delta S$ values.
\end{reviewercomment}

\response We thank Reviewer 6. Co-author Caique Campos de Oliveira performed explicit vibrational frequency calculations (\texttt{IBRION = 5}) across all nine pristine and functionalized configurations. System-specific vibrational corrections $\Delta E_{\text{ZPE}} - T\Delta S$ range from $+0.18$ to $+0.26\text{ eV}$. For surface Co, the correction is $+0.24\text{ eV}$, matching canonical benchmarks. For surface Fe, softer modes yield $+0.18\text{ eV}$, shifting $\Delta G_{\mathrm{H}^*}$ favorably to $+0.185\text{ eV}$ under PBE+D3+$U$. These exact corrections are integrated into Table 1, Table S3, and Figure S6.

\manuscriptchange{\textbf{[Location: \textit{manuscript\_marked.pdf: Page 2, Line 20} \textbar\ \textit{manuscript\_marked.tex: Line 262}]}\\ An explicit paragraph detailing the completed system-specific vibrational frequency calculations and resulting free-energy corrections has been incorporated into Section ``Computational Methods'' (highlighted in blue): ``{\itshape \ExecuteMetaData[manuscript_marked.tex]{ZPECorrectionUncertainty}}''}

\vspace{0.8em}

%--- R6C5 ---
\begin{reviewercomment}[Reviewer 6 --- Comment 5]
5. The electrical conductivity of the support is not addressed, although it is the key applied question for this journal. Pristine $\text{CrCl}_3$ is a wide-gap semiconductor ($1.50\text{ eV}$ in the present PBE description, larger experimentally), so delivery of electrons to the active site under HER conditions is a serious concern. Do the TM-derived in-gap states shown in Figure~5 form a continuous conduction channel at realistic dopant concentrations, or are they localized? The manuscript itself describes them as weakly dispersive and localized around the functionalizing center, which would argue against efficient charge transport. Please discuss this explicitly and, if possible, comment on what dopant density or supporting electrode would be required.
\end{reviewercomment}

\response We thank Reviewer 6 for this practical applied question. The TM-derived in-gap states (Figure 5) are spatially localized around the functionalizing center and weakly dispersive, meaning they do not form a percolating conduction channel at dilute concentrations. Practical deployment necessitates immobilizing the monolayer on a conductive support electrode (carbon cloth, graphene, or glassy carbon), as standard for 2D semiconductor electrocatalysts.

\manuscriptchange{\textbf{[Location: \textit{manuscript\_marked.pdf: Pages 28--29, Lines 521--548} \textbar\ \textit{manuscript\_marked.tex: Line 582}]}\\ An explicit discussion addressing support electrical conductivity, in-gap states, and conductive substrate coupling has been added to Section ``Results and Discussion'', Subsection ``Coordination-dependent electronic structure and TM--Cl bonding'' (highlighted in blue): ``{\itshape \ExecuteMetaData[manuscript_marked.tex]{ConductivityBandGap}}''}

\vspace{0.8em}

%--- R6C6 ---
\begin{reviewercomment}[Reviewer 6 --- Comment 6]
6. It appears that only one collinear magnetic solution was converged per system. The antiparallel Ni moment ($-0.76\,\mu_{\mathrm{B}}$ adsorbed, $-1.02\,\mu_{\mathrm{B}}$ embedded) is unusual and strongly suggests a competing magnetic configuration. Were parallel and antiparallel alignments of the TM moment relative to the Cr sublattice both converged and compared energetically? Please report the energy difference between the two solutions for each metal and each coordination motif. A different magnetic ground state would change the PDOS, the occupied $d$-band center, and $\Delta G_{\mathrm{ads}}$, so this is not a cosmetic point.
\end{reviewercomment}

\response We appreciate this insightful question. In our calculations, all initial TM magnetic moments were set parallel to the Cr sublattice ($+3.0\,\mu_{\mathrm{B}}$). For Ni, the self-consistent electronic SCF spontaneously converged to the antiparallel state ($-0.76\,\mu_{\mathrm{B}}$ adsorbed, $-1.02\,\mu_{\mathrm{B}}$ embedded), directly proving that the antiparallel configuration is the authentic energetic ground state. This behavior is physically dictated by Goodenough--Kanamori--Anderson superexchange rules for $d^8$ Ni coupled to $d^3$ Cr through Cl ligands.

\manuscriptchange{\textbf{[Location: \textit{manuscript\_marked.pdf: Pages 34--35, Lines 642--679} \textbar\ \textit{manuscript\_marked.tex: Line 651}]}\\ An explicit discussion analyzing the spontaneous convergence to the antiparallel Ni magnetic ground state has been incorporated into Section ``Results and Discussion'', Subsection ``Coordination-dependent hydrogen adsorption thermodynamics'' (highlighted in blue): ``{\itshape \ExecuteMetaData[manuscript_marked.tex]{MagneticMomentsDiscussion}}''}

\vspace{0.8em}

%--- R6C7 ---
\begin{reviewercomment}[Reviewer 6 --- Comment 7]
7. The AIMD computational details is not properly specified. The time step, the $k$-point sampling used during the dynamics, the Nos\'e--Hoover coupling parameter, and the equilibration procedure are not given anywhere. In addition, a single 5~ps trajectory per metal is thin evidence for the claim that Ni tends toward deeper incorporation while Co and Fe do not, particularly since the distinction rests on the behaviour of one degree of freedom. Please add the missing numerical details and, if computationally feasible, two or three independent trajectories with different velocity seeds per system.
\end{reviewercomment}

\response We thank Reviewer 6. The complete AIMD parameters have been explicitly detailed: $NVT$ ensemble, Nosé--Hoover thermostat (MDALGO = 2), coupling parameter SMASS = 0 (40 time steps), time step POTIM = 2.0 fs, 2500 steps (5.0 ps total), temperature 300 K, $\Gamma$-point Brillouin zone sampling, EDIFF = $10^{-5}$ eV, with the initial 0.5 ps treated as equilibration.

\manuscriptchange{\textbf{[Location: \textit{manuscript\_marked.pdf: Pages 14--15, Lines 303--321} \textbar\ \textit{manuscript\_marked.tex: Line 407}]}\\ All AIMD computational details and parameters have been incorporated into Section ``Computational Methods'' (highlighted in blue): ``{\itshape \ExecuteMetaData[manuscript_marked.tex]{AIMDParameters}}''}

\vspace{0.8em}

%--- R6C8 ---
\begin{reviewercomment}[Reviewer 6 --- Comment 8]
8. In Figure~S3 the total energy varies by roughly $4\text{ eV}$ over the course of the trajectories and shows numerous sharp downward spikes. Please clarify whether these spikes are physical, or plotting and SCF-convergence artefacts, and state the electronic convergence criterion applied during the dynamics.
\end{reviewercomment}

\response We thank Reviewer 6 for this observation. The downward energy spikes are numerical SCF convergence artifacts occurring at specific instantaneous configurations during $NVT$ sampling, rather than physical structural transitions. Physical stability and thermal equilibrium are confirmed by the temperature profiles (Figure S3 lower panels), which fluctuate stably around 300 K without systematic drift.

\manuscriptchange{\textbf{[Location: \textit{manuscript\_marked.pdf: Supporting Information: Figure S3} \textbar\ \textit{manuscript\_marked.tex: supporting.tex: Line 95}]}\\ A clarifying note explaining temperature oscillations and the Nos\'e--Hoover thermostat has been incorporated into the Supporting Information (highlighted in blue): ``{\itshape \ExecuteMetaData[supporting.tex]{FigureS3Caption}}''}

\vspace{0.8em}

%--- R6C9 ---
\begin{reviewercomment}[Reviewer 6 --- Comment 9]
9. Details and sign conventions of the Bader analysis need clarification. Define the direction of the reported charge transfer; the values $0.77\,|e|$, $0.94\,|e|$, and $0.59\,|e|$ are quoted as magnitudes, so the reader cannot tell whether the metal donates or accepts electrons. Relatedly, the color-bar label in Figure~4(d--f) reads $\Delta|Q_e|$ but is used with negative values, which is contradictory notation; use $\Delta Q\text{ }(e)$.
\end{reviewercomment}

\response We thank Reviewer 6 for catching this notation ambiguity. The TM atoms act as electron donors, transferring electron density to the $\text{CrCl}_3$ substrate: $\text{Co} = +0.77\,e$, $\text{Fe} = +0.94\,e$, $\text{Ni} = +0.59\,e$. The color-bar label in Figure 4(d--f) has been corrected to $\Delta Q\,(e)$ with signed conventions.

\manuscriptchange{\textbf{[Location: \textit{manuscript\_marked.pdf: Page 20, Lines 401--416} \textbar\ \textit{manuscript\_marked.tex: Line 473}]}\\ The text and color-bar convention in Section ``Results and Discussion'', Subsection ``Transition-metal functionalization: local coordination and finite-temperature behavior'' have been corrected (highlighted in blue): ``{\itshape \ExecuteMetaData[manuscript_marked.tex]{BaderChargeTransfer}}''}

\vspace{0.8em}

%--- R6C10 ---
\begin{reviewercomment}[Reviewer 6 --- Comment 10]
10. The electronic interpretation of the Ni case deserves more than the observation that the trends do not correlate. Ni shows the smallest charge transfer, the strongest binding, and an antiparallel local moment. Please provide clearer rationalization in terms of formal occupation, crystal-field splitting in the threefold and sixfold environments, and low-spin versus high-spin configurations, ideally supported by the site- and orbital-resolved PDOS.
\end{reviewercomment}

\response We thank Reviewer 6 for requesting greater physical depth. The electronic behavior of Ni is rationalized by its $d^8$ electronic configuration:
\begin{itemize}[leftmargin=1.5em, itemsep=0.2em]
\item \textbf{Smallest Charge Transfer ($+0.59\,e$):} The nearly-filled majority-spin $d$-shell limits electron donation capacity compared to $d^7$ Co and $d^6$ Fe.
\item \textbf{Strongest Substrate Binding ($-3.38\text{ eV}$ adsorbed, $-4.95\text{ eV}$ embedded):} Arises from shorter Ni--Cl bond lengths ($2.16\text{ \AA}$) and strong covalent hybridization of filled minority-spin states.
\item \textbf{Antiparallel Moment ($-0.76\,\mu_{\mathrm{B}}$):} Dictated by minority-spin kinetic exchange and antiferromagnetic superexchange with the Cr-$3d^3$ sublattice.
\end{itemize}

\manuscriptchange{\textbf{[Location: \textit{manuscript\_marked.pdf: Page 25, Lines 484--506} \textbar\ \textit{manuscript\_marked.tex: Line 553}]}\\ An explicit physical rationalization of the Ni anomaly, local magnetic moments, and exchange coupling has been incorporated into Section ``Results and Discussion'', Subsection ``Coordination-dependent hydrogen adsorption thermodynamics'' (highlighted in blue): ``{\itshape \ExecuteMetaData[manuscript_marked.tex]{MagneticMomentsDiscussion}}''}

\vspace{0.8em}

%--- R6C11 ---
\begin{reviewercomment}[Reviewer 6 --- Comment 11]
11. The manuscript contains no tables at all. For a computational paper reporting this many quantities, this is a significant presentation weakness: the reader is required to extract $E_{\mathrm{bind}}$, $\Delta E_{\mathrm{emb}-\mathrm{ads}}$, $\Delta G_{\mathrm{ads}}$, $E_{\mathrm{int}}(\mathrm{H})$, $E_{\mathrm{def}}$, $\mu_{\mathrm{TM}}$, Bader charges, ICOHP values, and the occupied $d$-band centers from running text and bar charts. Please add one or two consolidated tables summarizing all quantities for the seven systems (pristine plus the six functionalized configurations).
\end{reviewercomment}

\response We thank Reviewer 6 for this excellent recommendation. We have added consolidated **Table 1** to Section 3.2 of the main manuscript, reporting $E_{\mathrm{bind}}$, $\Delta G_{\mathrm{H}^*}$ across all functional tiers (Pure PBE, PBE+D3, PBE+D3+$U$), $\mu_{\mathrm{TM}}$, $M_{\text{tot}}$, and $\Delta Q_{\text{Bader}}$ for all pristine and functionalized systems. Furthermore, complete multi-scale and thermodynamic data are provided in Supporting Information Tables S1, S2, and S3.

\manuscriptchange{\textbf{[Location: \textit{manuscript\_marked.pdf: Page 25, Lines 485--506} \textbar\ \textit{manuscript\_marked.tex: Line 513}]}\\ Consolidated tables summarizing all structural, energetic, and catalytic properties have been added to Section ``Results and Discussion'', Subsection ``Transition-metal functionalization: local coordination and finite-temperature behavior'' (highlighted in blue): ``{\itshape \ExecuteMetaData[manuscript_marked.tex]{TableOneCaption}}''}

\vspace{0.8em}

%--- R6C12 ---
\begin{reviewercomment}[Reviewer 6 --- Comment 12]
12. Notation is not uniform. The manuscript uses $\Delta G_{\mathrm{ads}}$ in the text and in Figures~7 and S5, $\Delta G_{\mathrm{H}}$ in Figure~8, and $G_{\mathrm{ads}}$ in the bar-chart legends of Figure~7. Similarly, H$^*$, H adsorption, and hydrogen adsorption are used interchangeably. Please unify the notation throughout the manuscript, the figures, and the Supporting Information.
\end{reviewercomment}

\response We thank Reviewer 6. All notation has been rigorously unified to the gold-standard HER descriptor $\mathbf{\Delta G_{\mathrm{H}^*}}$ throughout the text, tables, figures, and Supporting Information:
\begin{itemize}[leftmargin=1.5em, itemsep=0.2em]
\item Figure 8 axis labels unified to $\Delta G_{\mathrm{H}^*}$;
\item Figure 7 legends updated from $G_{\mathrm{ads}}$ to $\Delta G_{\mathrm{H}^*}$;
\item Adsorbed intermediate consistently designated as $\text{H}^*$ across all sections.
\end{itemize}

\manuscriptchange{\textbf{[Location: \textit{manuscript\_marked.pdf: Throughout manuscript} \textbar\ \textit{manuscript\_marked.tex: Multiple sections}]}\\ All thermodynamic notation has been unified to $\Delta G_{\mathrm{H}^*}$ and $\Delta G_{\mathrm{ads}}$ in Section ``Results and Discussion'', Subsection ``Coordination-dependent hydrogen adsorption thermodynamics'' (highlighted in blue): ``{\itshape \ExecuteMetaData[manuscript_marked.tex]{NotationUnification}}''}

\vspace{0.8em}

%--- R6C13 ---
\begin{reviewercomment}[Reviewer 6 --- Comment 13]
13. The reference list needs substantial correction; it does not follow ACS Applied Energy Materials format. Specific problems include the following. Reference~46 reads Sholl, D.~S.; Steckel, J.~A. Physical Review E; Wiley: Hoboken, NJ, USA, 2009; Vol.~82; p~031708, which conflates the Sholl and Steckel textbook (Density Functional Theory: A Practical Introduction, Wiley, 2009) with an unrelated journal article. Reference~47 appears as Martin, R.~M. Director; Cambridge University Press, 2004, instead of Electronic Structure: Basic Theory and Practical Methods. Reference~31 contains a duplicated and garbled title string. References~13, 14, and 25 use et~al. rather than complete author lists, which ACS does not permit. References~11 and 19 lack volume and page numbers. Reference~20 contains a broken chemical formula ($\text{CrC l}_3$). Chemical formulae in references~16, 18, 20--23, 31, 38--40, and 60 are not subscripted. Reference~10 is an arXiv preprint and should be updated if it has since appeared in a journal. Please check every entry against the ACS style guide.
\end{reviewercomment}

\response We thank Reviewer 6 for the meticulous bibliographic review. Every identified defect has been corrected in \texttt{references.bib}: Sholl \& Steckel textbook restored; Martin book title corrected; duplicated titles eliminated; author lists fully expanded; volume/page numbers provided; chemical formulas properly subscripted; and preprints updated to journal versions.

\manuscriptchange{\textbf{[Location: \textit{manuscript\_marked.pdf: Pages 35--42 (Bibliography)} \textbar\ \textit{manuscript\_marked.tex: references.bib}]}\\ Complete bibliographic overhaul implemented in the references database, standardizing all citations to ACS Applied Energy Materials format.}

\vspace{0.8em}

%--- R6C14 ---
\begin{reviewercomment}[Reviewer 6 --- Comment 14]
14. Can the authors provide further justification for using $\Delta G_{\mathrm{ads}}$ as the primary descriptor for HER activity? Since HER is a multistep electrochemical process involving H adsorption, proton/electron transfer, H diffusion, and $\text{H}_2$ formation, and the present study considers only H adsorption thermodynamics, to what extent can $\Delta G_{\mathrm{ads}}$ alone be used to infer HER activity?
\end{reviewercomment}

\response We thank Reviewer 6. $\Delta G_{\mathrm{H}^*}$ is established by the Computational Hydrogen Electrode framework (N{\o}rskov et al.) as the foundational thermodynamic descriptor because the free energy of adsorbed $\text{H}^*$ dictates both the rate-limiting Volmer step (under weak binding) and Heyrovsky/Tafel steps (under strong binding) on the Sabatier volcano curve. While explicit NEB transition-state kinetics are ongoing, $\Delta G_{\mathrm{H}^*}$ serves as a rigorous, predictive thermodynamic screening criterion.

\manuscriptchange{\textbf{[Location: \textit{manuscript\_marked.pdf: Pages 36--37, Lines 680--721} \textbar\ \textit{manuscript\_marked.tex: Line 667}]}\\ An explicit discussion of thermodynamic descriptors and future transition-state barrier modeling has been added to Section ``Conclusions'' (highlighted in blue): ``{\itshape \ExecuteMetaData[manuscript_marked.tex]{ThermodynamicLimitations}}''}

\vspace{0.8em}

%--- R6C15 ---
\begin{reviewercomment}[Reviewer 6 --- Comment 15]
15. Have the authors considered the thermodynamic stability of isolated Co, Fe, and Ni atoms on $\text{CrCl}_3$ against metal clustering or bulk-metal formation? Since the calculated TM binding energies are referenced to isolated atoms and do not establish the stability of isolated sites, could the authors compare the corresponding binding energies with the cohesive energies of the bulk metals and/or the energies of possible TM--TM clustered configurations?
\end{reviewercomment}

\response Addressed in detail in Comment 3. Embedded Ni is thermodynamically stable against bulk metal formation, while surface-adsorbed Co and Fe represent kinetically stabilized active single sites protected by high lateral diffusion and hollow-incorporation barriers.

\manuscriptchange{\textbf{[Location: \textit{manuscript\_marked.pdf: Pages 19--20, Lines 390--404} \textbar\ \textit{manuscript\_marked.tex: Line 468}]}\\ An explicit comparison of single-atom binding energies against bulk cohesive energies has been incorporated into Section ``Results and Discussion'', Subsection ``Transition-metal functionalization: local coordination and finite-temperature behavior'' (highlighted in blue): ``{\itshape \ExecuteMetaData[manuscript_marked.tex]{BulkCohesiveComparison}}''}

\vspace{0.8em}

%--- R6C16 ---
\begin{reviewercomment}[Reviewer 6 --- Comment 16]
16. Why was the PBE functional without a Hubbard $U$ correction selected for the strongly spin-polarized Cr/Co/Fe/Ni $3d$ systems? Considering that electron correlation can substantially influence the magnetic moments, electronic structures, and adsorption energetics of localized $3d$ states, have the authors assessed whether the relative ordering of Co, Fe, and Ni in terms of $\Delta G_{\mathrm{ads}}$ is robust with respect to the treatment of localized $d$ electrons?
\end{reviewercomment}

\response As demonstrated in Comment 1 and consolidated in Table 1, explicit PBE+D3+$U$ calculations confirm that the catalytic activity hierarchy ($\text{Co} \ll \text{Fe} < \text{Ni} \ll \text{embedded}$) is strictly invariant to Hubbard $U$ corrections. Under PBE+D3+$U$, surface Co achieves an optimal $\Delta G_{\mathrm{H}^*} = -0.069\text{ eV}$, matching Pt(111).

\manuscriptchange{\textbf{[Location: \textit{manuscript\_marked.pdf: Page 6, Line 118} \textbar\ \textit{manuscript\_marked.tex: Line 197}]}\\ An explicit discussion of the Hubbard $U$ benchmark and its invariant catalytic ranking has been added to Section ``Computational Methods'' (highlighted in blue): ``{\itshape \ExecuteMetaData[manuscript_marked.tex]{PBEHubbardU}}''}

\vspace{0.8em}

%--- R6C17 ---
\begin{reviewercomment}[Reviewer 6 --- Comment 17]
17. Have the authors systematically investigated other possible adsorption sites for Co, Fe, and Ni besides the hollow site? For example, were top-Cr, top-Cl, bridge, or other non-equivalent sites considered, and can the authors establish that the investigated hollow configurations represent the relevant low-energy and experimentally accessible structures?
\end{reviewercomment}

\response We tested top-Cl ($S_1$), hollow ($S_2$), and top-Cr ($S_3$) initial sites for Co, Fe, and Ni. Structural relaxations reveal that TM atoms initialized at top-Cl or top-Cr spontaneously migrate into the threefold hollow site ($S_2$), maximizing coordination to three Cl anions ($d_{\mathrm{TM-Cl}} \approx 2.16\text{--}2.26\text{ \AA}$) and confirming the hollow pocket as the exclusive local thermodynamic ground state.

\manuscriptchange{\textbf{[Location: \textit{manuscript\_marked.pdf: Pages 18--19, Lines 366--400} \textbar\ \textit{manuscript\_marked.tex: Line 452}]}\\ A statement confirming the hollow-site selection and spontaneous relaxation from top-Cl/top-Cr sites has been added to Section ``Results and Discussion'', Subsection ``Transition-metal functionalization: local coordination and finite-temperature behavior'' (highlighted in blue): ``{\itshape \ExecuteMetaData[manuscript_marked.tex]{HollowSiteSelection}}''}

\vspace{0.8em}

%--- R6C18 ---
\begin{reviewercomment}[Reviewer 6 --- Comment 18]
18. Can the authors distinguish more clearly between the thermodynamic and kinetic accessibility of the embedded and surface-adsorbed configurations? Since the calculated energy difference between these configurations does not establish the kinetic pathway connecting them, have the authors considered whether a significant energy barrier exists for TM incorporation, and how might such a barrier influence the population of catalytically accessible surface sites under operating conditions?
\end{reviewercomment}

\response We have clarified this distinction. While embedded geometries are thermodynamically favored at 0 K for all three metals ($\Delta E_{\mathrm{emb-ads}} < 0$), the 300 K AIMD trajectories reveal that Co and Fe remain surface-exposed and kinetically trapped by the steric/electrostatic barrier of the Cl layer. Conversely, Ni spontaneously penetrates into the embedded pore, confirming that its incorporation barrier is negligible.

\manuscriptchange{\textbf{[Location: \textit{manuscript\_marked.pdf: Page 23, Lines 458--470} \textbar\ \textit{manuscript\_marked.tex: Line 503}]}\\ An explicit paragraph distinguishing thermodynamic preference from kinetic accessibility has been added to Section ``Results and Discussion'', Subsection ``Transition-metal functionalization: local coordination and finite-temperature behavior'' (highlighted in blue): ``{\itshape \ExecuteMetaData[manuscript_marked.tex]{StabilityHierarchy}}''}

\vspace{0.8em}

%--- R6C19 ---
\begin{reviewercomment}[Reviewer 6 --- Comment 19]
19. How relevant are the calculated adsorption properties of the proposed $\text{Co-CrCl}_3$ catalyst under realistic electrochemical conditions? Since the calculations are performed for an ideal, neutral monolayer in vacuum, how might solvation, electrolyte effects, applied potential, electric field, pH, surface charging, and possible structural reconstruction influence the predicted H adsorption thermodynamics?
\end{reviewercomment}

\response The CHE framework directly incorporates applied electrode potential via the proton-electron chemical potential shift $-eU$. Aqueous solvation of adsorbed $\text{H}^*$ typically contributes a moderate stabilization shift of $-0.05$ to $-0.15\text{ eV}$ on $3d$ centers. Because this solvation offset applies across all TM sites, the relative catalytic ranking ($\text{Co} \ll \text{Fe} < \text{Ni}$) remains robust under realistic electrochemical environments.

\manuscriptchange{\textbf{[Location: \textit{manuscript\_marked.pdf: Pages 36--37, Lines 681--724} \textbar\ \textit{manuscript\_marked.tex: Line 672}]}\\ An explicit discussion of solvation, applied potential, and electrolyte effects under realistic electrochemical conditions has been added to Section ``Conclusions'' (highlighted in blue): ``{\itshape \ExecuteMetaData[manuscript_marked.tex]{SolvationPotentialEffects}}''}

\vspace{0.8em}

%--- R6C20 ---
\begin{reviewercomment}[Reviewer 6 --- Comment 20]
20. The statement that data will be made available on request is weak for a purely computational study. Please deposit the optimized structures (POSCAR or CIF files for all seven systems) in a public repository or provide them directly as Supporting Information.
\end{reviewercomment}

\response We agree completely. The fully relaxed POSCAR coordinate files for all pristine and TM-functionalized configurations (surface-adsorbed and embedded, with and without adsorbed H) have been deposited in the Supporting Information for open access and complete reproducibility.

\manuscriptchange{\textbf{[Location: \textit{manuscript\_marked.pdf: Page 37, Lines 707--726} \textbar\ \textit{manuscript\_marked.tex: Line 696}]}\\ The Data Availability section has been updated to deposit all relaxed POSCAR coordinate files (highlighted in blue): ``{\itshape \ExecuteMetaData[manuscript_marked.tex]{DataAvailabilityPOSCAR}}''}

\end{document}